% concatenated from GODELSmImplication.tex

\documentclass[11pt]{article}
\usepackage{amsmath,amssymb,amsthm}
\usepackage{geometry}
\usepackage{stmaryrd}
 \usepackage{longtable}
\usepackage{graphicx}
\usepackage{caption}
\usepackage{booktabs}
\usepackage{hyperref}
\usepackage{url}
\usepackage{placeins}
\usepackage{float}
\geometry{margin=1in}

\newtheorem{definition}{Definition}
\newtheorem{proposition}{Proposition}
\newtheorem{corollary}{Corollary}
\newtheorem{theorem}{Theorem}
\newtheorem{lemma}{Lemma}
\newtheorem{remark}{Remark}
\renewcommand{\qedsymbol}{$\star$}
\newcommand{\ptop}{p_{\top}}
\newcommand{\pbot}{p_{\bot}}


\begin{document}
\begin{titlepage}

    \centering
    \vspace*{2.5cm}

    {\Large\bfseries Gödel Implication on Finite Chains: \\ Truth Tables and Catalan-Bracketing Enumerations}
    \vspace{0.5cm}
 
    {\Large\scshape Volkan Yildiz\par}
    \vspace{0.3cm}
    {\large\texttt{volkan@hotmail.co.uk}\par}
    \vfill
    {\normalsize \today \par}

\begin{abstract}
Fully bracketed implication terms on $n$ variables are evaluated in G\"odel $m$-valued logic on a finite chain, and we enumerate truth-table rows by output value across all Catalan bracketings. Using the Catalan decomposition, we derive a finite system of generating functions for these value counts and introduce a root-split refinement that records the ordered pair of truth values at the top implication, yielding $m^2$ pair classes. We prove that the associated generating functions share a common dominant square-root singularity, which implies a universal $n^{-3/2}$ asymptotic form with exponential growth rate $(4m)^n$ and a limiting output distribution as $n\to\infty$. The root-split refinement yields matching uniform asymptotics for the pair classes and gives a transparent factorization of the original counts.

\vfill
\noindent\textbf{Keywords:} Gödel many-valued logic, Gödel implication, truth tables, fully bracketed implications,
Catalan bracketings, generating functions, analytic combinatorics, singularity analysis,
asymptotic enumeration, limit laws, limit proportions

\medskip
\noindent\textbf{AMS 2020 Mathematics Subject Classification:}\\
Primary: 03B50, 05A15, 05A16, 05A19.\\
Secondary 05A10.
\end{abstract}




\end{titlepage}



\makeatletter
\newcommand{\oneLinePart}[1]{
  \cleardoublepage
  \refstepcounter{part}
  \thispagestyle{plain}
  \addcontentsline{toc}{part}{\partname~\thepart\quad #1}
  \begin{center}
    \makebox[\textwidth][c]{\Large\bfseries \partname~\thepart\quad #1}
  \end{center}
  \vspace{1.2em}
}
\makeatother




\newpage
 

\setcounter{tocdepth}{1} 
\tableofcontents
\bigskip

\newpage
\section{Introduction}



The combinatorial study of truth tables of propositional formulae connected by implication has been considered in the literature. In the classical two-valued setting, the
enumeration of propositional functions and the relationship between functions and
formulae are well understood. More recently, our attention has shifted to counting problems for families of formulae whose syntactic structure is restricted, while their semantic behaviour is analysed in detail.

In a series of papers \cite{CameronYildiz2010},  \cite{Yildiz2011b}, \cite{Yildiz2012}, we investigated the truth tables arising from all
fully bracketed implications on distinct variables, namely all bracketings of the word
\[
p_1 \Rightarrow p_2 \Rightarrow \cdots \Rightarrow p_n .
\]
The number of such bracketings is the Catalan number $C_{n-1}$, and the Catalan
decomposition provides a natural recursive framework for analysing the associated truth
tables. In the classical case and in Kleene's three-valued logic 
\cite{Yildiz2020}, this approach
yields explicit recurrences and algebraic generating functions for the numbers of truth
table entries of each output type, together with sharp asymptotics.


The aim of the present paper is to extend this programme to Gödel implication on a finite
chain. In contrast with the Kleene setting, Gödel logic is order--based: truth values form a
totally ordered set, and Gödel implication is determined by this order. This order--theoretic feature introduces a pronounced hierarchical behaviour in the output values and leads to a different combinatorial structure from the previously studied cases.
We begin with the four--valued Gödel chain
\[
V_4=\{0<a<b<1\},
\]
which serves as a toy model illustrating the new phenomena. For each \(n\), we consider all
fully bracketed implications on \(n\) distinct variables, evaluate them under all \(4^n\)
valuations into \(V_4\), and count the total numbers of occurrences of each truth value in the
output columns of the resulting truth tables. The Catalan decomposition again gives
recurrences; however, instead of decoupled scalar relations, the Gödel semantics naturally
lead to a stratified system reflecting the order \(0<a<b<1\). Using an explicit
\(\Phi/\Psi\)--iteration mechanism, we reduce this system to a single governing generating
function, obtain exact closed forms for the limiting proportions of each truth value (in nested
radicals), and provide numerical convergence tables supporting the asymptotic regime.

We extend the method to Gödel $m$--valued implication on the finite chain
$V_m=\{0=v_0<\cdots<v_{m-1}=1\}$.  The same $\Phi/\Psi$ calculus leads to a finite--depth
iteration of length $m-1$ that determines the full $m$--component output distribution.
In particular, we obtain explicit expressions for the limiting proportions $p_k(m)$ and show
that all output generating functions share a common square--root singularity at the dominant
point $r=1/(4m)$, yielding the universal coefficient decay $n^{-3/2}$.

We then analyse the large--$m$ limit.  The mass at the top value converges to
$\ptop(m)\to 1/\sqrt3$, while the bottom value satisfies $\pbot(m)\sim 1/m$.
More generally, fixed truth levels carry vanishing mass as $m\to\infty$; the correct scaling
is instead \emph{macroscopic}: cuts at rank $\lfloor t(m-1)\rfloor$ converge to a nontrivial
limit law with survival function $\Pr(T\ge t)=(1+2t)^{-1/2}$ for $0\le t<1$, together with an
atom of size $1/\sqrt3$ at $t=1$.

We conclude with a further refinement that records the ordered pair of root values
$\bigl(v(\varphi),v(\psi)\bigr)$.  This produces $m^2$ ``pair--count'' sequences (or
$\binom{m+1}{2}$ up to the symmetry $(i,j)\leftrightarrow(j,i)$), whose generating functions
are simple products of the level OGFs; they inherit the same dominant square--root singularity
and hence satisfy the same universal $n^{-3/2}$ transfer asymptotic.



\makeatletter
\@addtoreset{section}{part}
\@addtoreset{subsection}{part}
\makeatother

\renewcommand{\thesection}{\Roman{part}.\arabic{section}}
\renewcommand{\thesubsection}{\thesection.\arabic{subsection}}


\newpage

\oneLinePart{Gödel Four--Valued Implication}

\section{Recurrence Relations}
Let $p_1, p_2, \ldots ,p_n$ be distinct propositional variables.
Let \(V_4=\{0<a<b<1\}\) and consider the formula \(p_1 \Rightarrow_G p_2\).
Gödel implication is defined by
\[
x \Rightarrow_G y =
\begin{cases}
1, & \text{if } x \le y,\\
y, & \text{if } x > y .
\end{cases}
\]

\medskip
The induced operation table is shown in Table~\ref{tab:godel-op}.

\begin{table}[h]
\centering
\[
\begin{array}{c|cccc}
p_1 \Rightarrow_G p_2 & 0 & a & b & 1\\
\hline
0 & 1 & 1 & 1 & 1\\
a & 0 & 1 & 1 & 1\\
b & 0 & a & 1 & 1\\
1 & 0 & a & b & 1
\end{array}
\]
\caption{Gödel implication on \(V_4=\{0<a<b<1\}\).}
\label{tab:godel-op}
\end{table}

The full truth table, listing all valuations of \(p_1\) and \(p_2\), is given in the appendix;
see Table~\ref{tab:godel-n2}. Truth tables for all valuations of \(p_1\), \(p_2\), and \(p_3\)
are likewise provided in the appendix; see Table~\ref{tab:godel4-n3-sidebyside}.
\\\\
In the Gödel four--valued chain \(V_4=\{0<a<b<1\}\) we write
\[
f_n=\#\text{ of zeros},\qquad a_n=\#\text{ of \(a\)'s},\qquad b_n=\#\text{ of \(b\)'s},\qquad
t_n=\#\text{ of ones},
\]
counted across the output columns of the truth tables of all fully bracketed implications
on \(n\) variables. Thus
\[
g_n=f_n+a_n+b_n+t_n,
\qquad g_n=4^n C_{n-1}.
\]

By the Catalan decomposition of bracketings, we have
\[
g_n=\sum_{i=1}^{n-1} g_i g_{n-i}
   =\sum_{i=1}^{n-1}(f_i+a_i+b_i+t_i)(f_{n-i}+a_{n-i}+b_{n-i}+t_{n-i}).
\]

Using the Gödel implication
\[
x\Rightarrow_G y=
\begin{cases}
1,& x\le y,\\
y,& x>y,
\end{cases}
\]
all outputs other than 1 come only from the following product contributions.

\medskip

\noindent\textbf{Output $b$:}
\[
b_n=\sum_{i=1}^{n-1}
\underbrace{t_i\,b_{n-i}}_{\text{$1\Rightarrow b=b$}}.
\]

\noindent\textbf{Output $a$:}
\[
a_n=\sum_{i=1}^{n-1}\left(
\underbrace{b_i\,a_{n-i}}_{\text{$b\Rightarrow a=a$}}
+
\underbrace{t_i\,a_{n-i}}_{\text{$1\Rightarrow a=a$}}
\right)
=\sum_{i=1}^{n-1}\underbrace{(b_i+t_i)\,a_{n-i}}_{\text{output $a$}}.
\]

\noindent\textbf{Output $0$:}
\[
f_n=\sum_{i=1}^{n-1}\left(
\underbrace{a_i\,f_{n-i}}_{\text{$a\Rightarrow 0=0$}}
+
\underbrace{b_i\,f_{n-i}}_{\text{$b\Rightarrow 0=0$}}
+
\underbrace{t_i\,f_{n-i}}_{\text{$1\Rightarrow 0=0$}}
\right)
=\sum_{i=1}^{n-1}\underbrace{(a_i+b_i+t_i)\,f_{n-i}}_{\text{output $0$}}.
\]

\noindent\textbf{Output $1$:}
\[
t_n=g_n-f_n-a_n-b_n.
\]

Equivalently,
\[
t_n=\sum_{i=1}^{n-1}\Big(
g_i g_{n-i}
-
t_i b_{n-i}
-
(b_i+t_i)a_{n-i}
-
(a_i+b_i+t_i)f_{n-i}
\Big).
\]

\medskip

We define the ordinary generating functions
\[
F(x)=\sum_{n\ge1} f_n x^n,\qquad
A(x)=\sum_{n\ge1} a_n x^n,\qquad
B(x)=\sum_{n\ge1} b_n x^n,\qquad
T(x)=\sum_{n\ge1} t_n x^n,
\]
and
\[
G(x)=\sum_{n\ge1} g_n x^n.
\]

Since $g_n=4^n C_{n-1}$, we have
\[
G(x)=\sum_{n\ge 1}4^n C_{n-1}x^n=\frac{1-\sqrt{1-16x}}{2}.
\]

Our OGFs start at $n=1$ while the Catalan recurrences hold for $n\ge 2$, so the $n=1$
terms must be retained. For example,
$b_n=\sum_{i=1}^{n-1} t_i b_{n-i}$ gives $B(x)-x=T(x)B(x)$, i.e.\ $B(x)=x+T(x)B(x)$.
Similarly,
\[
A(x)=x+(B(x)+T(x))A(x),\qquad
F(x)=x+(A(x)+B(x)+T(x))F(x),
\]
and
\[
T(x)=G(x)-F(x)-A(x)-B(x).
\]


\subsection{Component sequences: the four--valued case}
Recall that
\[
g_n=f_n+a_n+b_n+t_n.
\]

The total sequence satisfies
\[
g_n=4^n C_{n-1}\sim \frac{16^n}{4\sqrt{\pi}\,n^{3/2}}.
\]
Numerical evaluation of the recurrences shows that each component grows on the
same exponential scale as \(g_n\), and the ratios \(f_n/g_n\), \(a_n/g_n\), \(b_n/g_n\),
\(t_n/g_n\) rapidly stabilise. For \(n=20\) one obtains approximately
\[
\frac{f_n}{g_n}\approx 0.1498,\qquad
\frac{a_n}{g_n}\approx 0.1006,\qquad
\frac{b_n}{g_n}\approx 0.0729,\qquad
\frac{t_n}{g_n}\approx 0.6767.
\]
This shows us that, in the four valued case, the value \(1\) dominates asymptotically,
while the three lower truth values occur with smaller but non--negligible limiting
frequencies.


\section{The $\Phi$--reduction and the governing generating function}

Consider the Gödel four--valued chain \(V_4=\{0<a<b<1\}\). Let
\[
F(x)=\sum_{n\ge1} f_n x^n,\qquad
A(x)=\sum_{n\ge1} a_n x^n,\qquad
B(x)=\sum_{n\ge1} b_n x^n,\qquad
T(x)=\sum_{n\ge1} t_n x^n
\]
denote the component generating functions corresponding to the output values
\(0,a,b,1\), respectively. Set
\[
G(x):=F(x)+A(x)+B(x)+T(x)
=\sum_{n\ge1} g_n x^n,
\qquad
g_n=4^n C_{n-1},
\]
so that
\begin{equation}\label{eq:G-closed}
G(x)=\sum_{n\ge1}4^nC_{n-1}x^n=\frac{1-\sqrt{1-16x}}{2}.
\end{equation}

\medskip

\noindent
We now re-introduce the functional equations (including the \(n=1\) terms).
The Gödel--\(4\) Catalan decompositions give:
\begin{equation}\label{eq:GF-system-godel4}
B(x)=x+T(x)B(x),\qquad
A(x)=x+(B(x)+T(x))A(x),\qquad
F(x)=x+(A(x)+B(x)+T(x))F(x),
\end{equation}
together with
\begin{equation}\label{eq:T-from-total}
T(x)=G(x)-F(x)-A(x)-B(x).
\end{equation}
\begin{proposition}[Recovering the remaining generating functions]\label{prop:recover-godel4}
Once \(T(x)\) is known, the remaining component generating functions are
\begin{equation}\label{eq:recover-ABF}
B(x)=\frac{x}{1-T(x)},\qquad
A(x)=\frac{x}{1-(B(x)+T(x))},\qquad
F(x)=\frac{x}{1-(A(x)+B(x)+T(x))}.
\end{equation}
\end{proposition}

\begin{proof}
These are immediate rearrangements of \eqref{eq:GF-system-godel4}.
\end{proof}

\begin{proposition}[$\Phi$--reduction: one equation in one unknown]\label{prop:phi-reduction-godel4}
Define the tail sums
\[
H_3(x):=T(x),\qquad
H_2(x):=B(x)+T(x),\qquad\]
\[
H_1(x):=A(x)+B(x)+T(x),\qquad
H_0(x):=F(x)+A(x)+B(x)+T(x)=G(x),
\]
and define the map
\begin{equation}\label{eq:Phi-def}
\Phi_x(u):=u+\frac{x}{1-u}.
\end{equation}
Then
\[
H_2(x)=\Phi_x(H_3(x)),\qquad
H_1(x)=\Phi_x(H_2(x)),\qquad
H_0(x)=\Phi_x(H_1(x)).
\]
Equivalently,
\begin{equation}\label{eq:G-Phi3-T}
\boxed{\,G(x)=\Phi_x^{\,3}\!\big(T(x)\big).\,}
\end{equation}
\end{proposition}

\begin{proof}
From \eqref{eq:GF-system-godel4} we obtain
\[
B(x)=\frac{x}{1-T(x)}.
\]
Hence
\[
H_2(x)=B(x)+T(x)=T(x)+\frac{x}{1-T(x)}=\Phi_x(H_3(x)).
\]
and, from \eqref{eq:GF-system-godel4},
\[
A(x)=\frac{x}{1-(B(x)+T(x))}=\frac{x}{1-H_2(x)},
\]
so
\[
H_1(x)=A(x)+H_2(x)=H_2(x)+\frac{x}{1-H_2(x)}=\Phi_x(H_2(x)).
\]
Finally,
\[
F(x)=\frac{x}{1-(A(x)+B(x)+T(x))}=\frac{x}{1-H_1(x)},
\]
so
\[
H_0(x)=F(x)+H_1(x)=H_1(x)+\frac{x}{1-H_1(x)}=\Phi_x(H_1(x)).
\]
Iterating gives \(H_0=\Phi_x^{3}(H_3)\), i.e. \eqref{eq:G-Phi3-T}.
\end{proof}

\medskip 

\begin{proposition}[Explicit inverse and nested radicals]\label{prop:Psi-inverse-godel4}
The map \(\Phi_x\) in \eqref{eq:Phi-def} is invertible on formal power series with zero
constant term. Its inverse branch is
\begin{equation}\label{eq:Psi-def}
\Psi_x(v):=\frac{1+v-\sqrt{(1-v)^2+4x}}{2},
\end{equation}
characterised by \(\Psi_x(0)=0\) and \(\Phi_x(\Psi_x(v))=v\).
Thus,
\begin{equation}\label{eq:T-Psi3-G}
\boxed{\,T(x)=\Psi_x^{\,3}\!\big(G(x)\big),\qquad
G(x)=\frac{1-\sqrt{1-16x}}{2}.\,}
\end{equation}
\end{proposition}

\begin{proof}
If \(v=\Phi_x(u)=u+\frac{x}{1-u}\), then \(v(1-u)=u(1-u)+x\), i.e.
\(u^2-(1+v)u+(v-x)=0\). Solving this quadratic gives
\[
u=\frac{1+v\pm\sqrt{(1-v)^2+4x}}{2}.
\]
The branch with \(u(0)=0\) is obtained by choosing the minus sign, giving \eqref{eq:Psi-def}.
Applying \(\Psi_x\) three times to \eqref{eq:G-Phi3-T} gives \eqref{eq:T-Psi3-G}.
\end{proof}

\medskip

%77777777777777
%77777777777777


\section{Asymptotics}

Let \(r=1/16\) be the dominant singularity of \(G(x)\). From \eqref{eq:G-closed} we have the
local expansion
\begin{equation}\label{eq:G-local}
G(x)=\frac12-\frac12\sqrt{1-\frac{x}{r}}+O\!\left(1-\frac{x}{r}\right)\qquad (x\to r^-).
\end{equation}

\newpage

\begin{theorem}[Limiting proportion of $1$'s and asymptotics]\label{thm:limit-ones-godel4}
Let \(\ptop:=\lim_{n\to\infty} \frac{t_n}{g_n}\), where \(t_n=[x^n]T(x)\) and
\(g_n=[x^n]G(x)=4^nC_{n-1}\). Then \(p\) exists and
\[
t_n \sim \frac{p}{4}\cdot \frac{16^n}{\sqrt{\pi}\,n^{3/2}}.
\]
\end{theorem}

\begin{proof}
By \eqref{eq:G-local} and the transfer theorem for square--root singularities
\cite[Ch.~VI]{FlajoletSedgewick2009},
\[
g_n=[x^n]G(x)\sim \frac{16^n}{4\sqrt{\pi}\,n^{3/2}}.
\]
Since \(T(x)=\Psi_x^{\,3}(G(x))\) by \eqref{eq:T-Psi3-G} and \(\Psi_x(v)\) is analytic in \(v\)
near the values encountered for \(x\) in a neighborhood of \(r\), the singular expansion
\eqref{eq:G-local} propagates through composition and yields
\[
T(x)=T(r)-\lambda \sqrt{1-\frac{x}{r}}+O\!\left(1-\frac{x}{r}\right)\qquad (x\to r^-)
\]
for some \(\lambda>0\). Applying the same transfer theorem gives
\[
t_n=[x^n]T(x)\sim \frac{\lambda}{\sqrt{\pi}}\frac{16^n}{n^{3/2}}.
\]
Taking the ratio shows \(\frac{t_n}{g_n}\to 4\lambda=:p\), and the stated form of \(t_n\)
follows.
\end{proof}

\subsection{Gödel--$4$: exact limiting proportions and asymptotic constants}

Recall
\[
G(x)=F(x)+A(x)+B(x)+T(x)=\frac{1-\sqrt{1-16x}}{2}.
\]
Set \(r=\frac1{16}\) and recall the inverse map
\[
\Psi_x(v)=\frac{1+v-\sqrt{(1-v)^2+4x}}{2},
\]
so that \(T(x)=\Psi_x^{\,3}(G(x))\) (Proposition~\ref{prop:Psi-inverse-godel4}).  At the
dominant point \(x=r\) we have \(4r=\frac14\) and thus
\begin{equation}\label{eq:Psi-at-r}
\Psi_r(v)=\frac{1+v-\sqrt{(1-v)^2+\frac14}}{2}.
\end{equation}

\medskip
\noindent\textbf{Auxiliary radicals.}
Define
\[
\beta:=\sqrt{7+2\sqrt2},\qquad
\alpha:=1+\sqrt2+\beta,\qquad
\gamma:=\sqrt{16+\alpha^2}.
\]

\begin{lemma}[Explicit evaluation of the iterates at \(x=r\)]\label{lem:iterates-at-r}
Let \(w_0:=G(r)=\frac12\) and \(w_{k+1}:=\Psi_r(w_k)\) for \(k=0,1,2\). Then
\[
w_0=\frac12,\qquad
w_1=\frac34-\frac{\sqrt2}{4},\qquad
w_2=\frac78-\frac{\sqrt2}{8}-\frac{\beta}{8},\qquad
w_3=\frac{15}{16}-\frac{\sqrt2}{16}-\frac{\beta}{16}-\frac{\gamma}{16},
\]
where \(w_3=T(r)\).
\end{lemma}

\begin{proof}
We compute successively using \eqref{eq:Psi-at-r}.

\smallskip\noindent
\emph{Step 1.} Since \(w_0=\frac12\),
\[
\sqrt{(1-w_0)^2+\frac14}=\sqrt{\left(\frac12\right)^2+\frac14}=\sqrt{\frac12}=\frac{\sqrt2}{2},
\]
hence
\[
w_1=\Psi_r\!\left(\frac12\right)
=\frac{1+\frac12-\frac{\sqrt2}{2}}{2}
=\frac34-\frac{\sqrt2}{4}.
\]

\smallskip\noindent
\emph{Step 2.} Here \(1-w_1=\frac{1+\sqrt2}{4}\), so
\[
(1-w_1)^2+\frac14=\frac{(1+\sqrt2)^2}{16}+\frac14=\frac{7+2\sqrt2}{16},
\qquad
\sqrt{(1-w_1)^2+\frac14}=\frac{\beta}{4}.
\]
Substituting into \eqref{eq:Psi-at-r} yields
\[
w_2=\Psi_r(w_1)=\frac{1+w_1-\frac{\beta}{4}}{2}
=\frac78-\frac{\sqrt2}{8}-\frac{\beta}{8}.
\]

\smallskip\noindent
\emph{Step 3.} Next \(1-w_2=\frac{1+\sqrt2+\beta}{8}=\frac{\alpha}{8}\), hence
\[
(1-w_2)^2+\frac14=\frac{\alpha^2}{64}+\frac14=\frac{\alpha^2+16}{64},
\qquad
\sqrt{(1-w_2)^2+\frac14}=\frac{\gamma}{8}.
\]
Substituting into \eqref{eq:Psi-at-r} gives
\[
w_3=\Psi_r(w_2)
=\frac{1+w_2-\frac{\gamma}{8}}{2}
=\frac{15}{16}-\frac{\sqrt2}{16}-\frac{\beta}{16}-\frac{\gamma}{16}.
\]
\end{proof}

\begin{proposition}[Derivative product and the constant \(\ptop \)]\label{prop:p1-product}
Differentiating \(\Psi_x(v)\) gives
\begin{equation}\label{eq:dPsi}
\partial_v\Psi_x(v)
=\frac12\left(1+\frac{1-v}{\sqrt{(1-v)^2+4x}}\right).
\end{equation}
Evaluating at \(x=r\) along \(w_0,w_1,w_2\) yields
\[
\partial_v\Psi_r(w_0)=\frac{2+\sqrt2}{4},\qquad
\partial_v\Psi_r(w_1)=\frac{\alpha}{2\beta},\qquad
\partial_v\Psi_r(w_2)=\frac{\alpha+\gamma}{2\gamma}.
\]
Consequently,
\begin{equation}\label{eq:p1-product}
\ptop :=\lim_{n\to\infty}\frac{t_n}{g_n}
=\prod_{k=0}^{2}\partial_v\Psi_r(w_k)
=\frac{(2+\sqrt2)\,\alpha(\alpha+\gamma)}{16\,\beta\,\gamma}.
\end{equation}
\end{proposition}
\begin{proof}
A direct differentiation and substitution using Lemma~\ref{lem:iterates-at-r}.
\end{proof}


\begin{theorem}[Exact limiting proportions for Gödel--$4$]\label{thm:godel4-exact-proportions}
The limits
\[
\pbot  :=\lim_{n\to\infty}\frac{f_n}{g_n},\qquad
p_a:=\lim_{n\to\infty}\frac{a_n}{g_n},\qquad
p_b:=\lim_{n\to\infty}\frac{b_n}{g_n},\qquad
\ptop :=\lim_{n\to\infty}\frac{t_n}{g_n}
\]
exist and satisfy \(\pbot +p_a+p_b+\ptop =1\). Moreover, they are given exactly by
\[
\pbot=\frac12-\frac{\sqrt2}{4},
\qquad
p_a=\frac{2+\sqrt2}{2\alpha\beta},
\qquad
p_b=\frac{(2+\sqrt2)\alpha}{\beta\,\gamma(\alpha+\gamma)},
\qquad
\ptop =\frac{(2+\sqrt2)\alpha(\alpha+\gamma)}{16\,\beta\,\gamma}.
\]
In particular,
\[
\ptop \approx 0.684122210733017786\ldots\, .
\]
\end{theorem}

\begin{proof}
The constant \(p_1\) is given by Proposition~\ref{prop:p1-product}.
For the remaining components, use the recovery relations
\[
B(x)=\frac{x}{1-T(x)},\qquad
A(x)=\frac{x}{1-(B(x)+T(x))},\qquad
F(x)=\frac{x}{1-(A(x)+B(x)+T(x))},
\]
together with the fact that \(T(x)\) has a square--root singularity at \(x=r\) with
\(T(r)\in(0,1)\). Each of \(B,A,F\) is therefore analytic in \(T\) near \(T(r)\) and inherits
the same dominant square--root singularity at \(x=r\). Applying the transfer theorem to
each component and taking coefficient ratios gives the existence of \(p_0,p_a,p_b\), and a
direct simplification yields the stated closed forms. Finally,
\(\pbot +p_a+p_b+\ptop =1\) follows from \(G=F+A+B+T\).
\end{proof}

\begin{corollary}[False rows and non--true rows]\label{cor:false-rows-godel4}
Define \(u_n:=f_n\) and \(s_n:=g_n-t_n=f_n+a_n+b_n\). Then
\[
\lim_{n\to\infty}\frac{u_n}{g_n}=\pbot,
\qquad
\lim_{n\to\infty}\frac{s_n}{g_n}=1-\ptop=\pbot+p_a+p_b,
\]
so the limiting proportion of non--true rows is
\[
1-\ptop \approx 0.315877789266982214\ldots\, .
\]
Equivalently,
\[
S(x):=\sum_{n\ge1}s_nx^n=G(x)-T(x).
\]
\end{corollary}

\begin{theorem}[Exact asymptotic constants]\label{thm:godel4-exact-asymptotics}
As \(n\to\infty\),
\[
g_n\sim \frac{16^n}{4\sqrt{\pi}\,n^{3/2}},
\]
and each component satisfies
\[
f_n\sim \frac{\pbot}{4}\,\frac{16^n}{\sqrt{\pi}\,n^{3/2}},\qquad
a_n\sim \frac{p_a}{4}\,\frac{16^n}{\sqrt{\pi}\,n^{3/2}},\qquad
b_n\sim \frac{p_b}{4}\,\frac{16^n}{\sqrt{\pi}\,n^{3/2}},\qquad
t_n\sim \frac{\ptop}{4}\,\frac{16^n}{\sqrt{\pi}\,n^{3/2}},
\]
where \(\pbot, p_a,p_b,\ptop \) are as in Theorem~\ref{thm:godel4-exact-proportions}.
\end{theorem}

\begin{proof}
The asymptotic for \(g_n\) follows from \eqref{eq:G-local} by the transfer theorem, \cite[Chapter.~VI]{FlajoletSedgewick2009}.
Each component has the same dominant square--root singularity at \(x=r\), and the limiting
ratios \(f_n/g_n\to \pbot \), \(a_n/g_n\to p_a\), \(b_n/g_n\to p_b\), \(t_n/g_n\to \ptop \) hold by
Theorem~\ref{thm:godel4-exact-proportions}. The stated component asymptotics follow
immediately.
\end{proof}


\newpage
\begin{table}[t]
\centering
\small
\renewcommand{\arraystretch}{1.05}
\begin{tabular}{rcccc}
\hline
\(n\) & \(f_n/g_n\) & \(a_n/g_n\) & \(b_n/g_n\) & \(t_n/g_n\) \\
\hline
10  & 0.153354 & 0.103186 & 0.074472 & 0.668988 \\
20  & 0.149830 & 0.100597 & 0.072869 & 0.676705 \\
30  & 0.148687 & 0.099760 & 0.072343 & 0.679210 \\
40  & 0.148121 & 0.099347 & 0.072082 & 0.680450 \\
50  & 0.147783 & 0.099100 & 0.071926 & 0.681191 \\
75  & 0.147335 & 0.098774 & 0.071718 & 0.682173 \\
100 & 0.147112 & 0.098611 & 0.071615 & 0.682662 \\
150 & 0.146890 & 0.098449 & 0.071511 & 0.683150 \\
200 & 0.146779 & 0.098368 & 0.071460 & 0.683394 \\
250 & 0.146712 & 0.098320 & 0.071429 & 0.683540 \\
\hline
\(\infty\) & 0.146447 & 0.098126 & 0.071305 & 0.684122 \\
\hline
\end{tabular}
\caption{Convergence of the empirical proportions \(f_n/g_n,a_n/g_n,b_n/g_n,t_n/g_n\) to the
limit (convergence) constants \(\pbot ,p_a,p_b,\ptop \). The limiting values are shown to six
decimal places (their exact nested--radical forms are given in
Theorem~\ref{thm:godel4-exact-proportions}).}
\label{tab:g4-proportions}
\end{table}



%\\\\\\\\\\\\\\\\\\\\\\\\\\\\\\\\\\\\\\\\\\\\\\\\\\\
%\\\\\\\\\\\\\\\\\\\\\\\\\\\\\\\\\\\\\\\\\\\\\\\\\\\


\oneLinePart{Gödel $m$--Valued Implication}
\label{sec:godel-m}

\section{ $\;$ Recurrences, generating functions, and the $\Phi$--calculus}

Fix an integer \(m\ge2\) and let
\[
V_m=\{v_0<v_1<\cdots<v_{m-1}\},\qquad v_0=0,\quad v_{m-1}=1,
\]
be the Gödel chain of \(m\) truth values. The Gödel implication \(\Rightarrow_G:V_m\times V_m\to V_m\) is
\[
v_p\Rightarrow_G v_q=
\begin{cases}
v_{m-1},& p\le q,\\
v_q,& p>q.
\end{cases}
\]

Let \(\mathcal{B}_n\) be the set of full bracketings of
\[
p_1\Rightarrow p_2\Rightarrow\cdots\Rightarrow p_n,
\]
so \(|\mathcal{B}_n|=C_{n-1}\). A valuation is a map \(\nu:\{p_1,\dots,p_n\}\to V_m\), hence there are \(m^n\) valuations.

For each \(j\in\{0,1,\dots,m-1\}\) and \(n\ge1\), let \(g_n^{(j)}\) denote the total number of occurrences of the value \(v_j\) in the output columns of the truth tables of all bracketings in \(\mathcal{B}_n\), taken over all valuations. Define the total
\[
g_n:=\sum_{j=0}^{m-1} g_n^{(j)}.
\]
Since there are \(C_{n-1}\) bracketings and \(m^n\) valuations, we have
\[
g_n=m^n C_{n-1}.
\]



\subsection*{Catalan recurrences}
For \(n\ge2\), each bracketing splits uniquely into a left part on \(i\) variables and a right part on \(n-i\) variables, \(1\le i\le n-1\). Consequently,
\[
g_n=\sum_{i=1}^{n-1} g_i g_{n-i}.
\]

For each \(j\le m-2\), the value \(v_j\) can arise only when the right subformula evaluates to \(v_j\) and the left subformula evaluates to a \emph{strictly larger} truth value. Hence, for \(n\ge2\) and \(j=0,1,\dots,m-2\),
\begin{equation}\label{eq:godel-m-rec}
g_n^{(j)}=\sum_{i=1}^{n-1}\Big(\sum_{p=j+1}^{m-1} g_i^{(p)}\Big)\, g_{n-i}^{(j)}.
\end{equation}
The top value is determined by complement:
\begin{equation}\label{eq:godel-m-top}
g_n^{(m-1)}=g_n-\sum_{j=0}^{m-2} g_n^{(j)}.
\end{equation}
Initial conditions are
\[
g_1^{(j)}=1\quad (j=0,1,\dots,m-1),\qquad g_1=m.
\]

\subsection*{Generating functions}
Define the component generating functions
\[
G_j(x):=\sum_{n\ge1} g_n^{(j)}x^n\qquad (j=0,1,\dots,m-1),
\]
and the total
\[
G(x):=\sum_{n\ge1} g_n x^n=\sum_{j=0}^{m-1}G_j(x).
\]
Since \(g_n=m^n C_{n-1}\), we have the closed form
\begin{equation}\label{eq:godel-m-G}
G(x)=\sum_{n\ge1}m^nC_{n-1}x^n=\frac{1-\sqrt{1-4mx}}{2}.
\end{equation}

Multiplying \eqref{eq:godel-m-rec} by \(x^n\) and summing over \(n\ge2\) gives, for each \(j=0,1,\dots,m-2\),
\begin{equation}\label{eq:godel-m-Gj}
G_j(x)=x+H_{j+1}(x)\,G_j(x),
\qquad
H_{j+1}(x):=\sum_{p=j+1}^{m-1}G_p(x).
\end{equation}
Equivalently,
\[
G_j(x)=\frac{x}{1-H_{j+1}(x)}\qquad (0\le j\le m-2),
\qquad
G_{m-1}(x)=G(x)-\sum_{j=0}^{m-2}G_j(x).
\]

\subsection*{The $\Phi$--calculus (finite depth \(m-1\))}
Define the tail sums
\[
H_k(x):=\sum_{j=k}^{m-1}G_j(x)\qquad (k=0,1,\dots,m-1),
\]
so that \(H_0(x)=G(x)\) and \(H_{m-1}(x)=G_{m-1}(x)\). Define also
We use the map $\Phi_x(u):=u+\frac{x}{1-u}$ defined earlier in \eqref{eq:Phi-def}.

\begin{proposition}[Tail--sum iteration]\label{prop:godel-m-phi}
For each \(k=0,1,\dots,m-2\),
\[
H_k(x)=\Phi_x\!\big(H_{k+1}(x)\big).
\]
In particular, writing \(U(x):=H_{m-1}(x)=G_{m-1}(x)\),
\[
\boxed{\,G(x)=\Phi_x^{\,m-1}\!\big(U(x)\big).\,}
\]
\end{proposition}

\begin{proof}
From \eqref{eq:godel-m-Gj} we have \((1-H_{k+1})G_k=x\), hence \(G_k=\frac{x}{1-H_{k+1}}\).
Then
\[
H_k=G_k+H_{k+1}=H_{k+1}+\frac{x}{1-H_{k+1}}=\Phi_x(H_{k+1}).
\]
Iterating from \(k=m-2\) down to \(k=0\) gives \(H_0=\Phi_x^{m-1}(H_{m-1})\), i.e.
\(G=\Phi_x^{m-1}(U)\).
\end{proof}

\begin{proposition}[Explicit inverse and nested radicals]\label{prop:godel-m-psi}
The map \(\Phi_x\) is invertible on formal power series with zero constant term. Its inverse
branch is
\[
\Psi_x(v):=\frac{1+v-\sqrt{(1-v)^2+4x}}{2},
\qquad \Psi_x(0)=0,
\]
so that \(\Phi_x(\Psi_x(v))=v\). Consequently,
\[
\boxed{\,U(x)=G_{m-1}(x)=\Psi_x^{\,m-1}\!\big(G(x)\big),\qquad
G(x)=\frac{1-\sqrt{1-4mx}}{2}.\,}
\]
Moreover, once \(U\) is known, all tail sums are recovered by
\[
H_{m-1}=U,\qquad H_k=\Phi_x(H_{k+1})\ (k=m-2,\dots,0),
\]
and the components by
\[
G_k=H_k-H_{k+1}\quad (k=0,\dots,m-2),\qquad G_{m-1}=H_{m-1}.
\]
\end{proposition}

\begin{proof}
The identity \(v=\Phi_x(u)=u+\frac{x}{1-u}\) is equivalent to
\(u^2-(1+v)u+(v-x)=0\), hence
\[
u=\frac{1+v\pm\sqrt{(1-v)^2+4x}}{2}.
\]
The choice \(u(0)=0\) forces the minus sign, giving \(u=\Psi_x(v)\). Applying \(\Psi_x\)
successively \(m-1\) times to \(G=\Phi_x^{m-1}(U)\) gives \(U=\Psi_x^{m-1}(G)\).
The recovery formulas follow from the definitions.
\end{proof}


\section{ $\;$Asymptotics}


In this section we investigate how the limiting fraction of outputs equal to \(1=v_{m-1}\)
depends on \(m\). Fix \(m\ge2\) and let \(v_0<\cdots<v_{m-1}=1\) be the Gödel chain.  Let
\(g_n=m^nC_{n-1}\) be the total number of truth-table entries across all Gödel--\(m\) truth
tables for fully bracketed implications on \(n\) variables.  For \(j\in\{0,1,\dots,m-1\}\)
define
\[
g_n^{(j)}:=\#\{\text{output entries equal to }v_j\},
\]
and call
\[
\ptop(m):=\lim_{n\to\infty}\frac{g_n^{(m-1)}}{g_n}
\]
the \emph{convergence constant} (limiting proportion) of outputs equal to \(1\).
Since
\[
g_n\sim \frac{(4m)^n}{4\sqrt{\pi}\,n^{3/2}}\qquad (n\to\infty),
\]
the corresponding asymptotic coefficient is \(c_1(m)=\ptop(m)/4\) in
\[
g_n^{(m-1)}\sim \frac{\ptop(m)}{4}\cdot \frac{(4m)^n}{\sqrt{\pi}\,n^{3/2}}
\qquad (n\to\infty).
\]

\medskip

\noindent
For exact computation of \(\ptop(m)\) via the \(\Psi\)--chain, we set \(r=\frac{1}{4m}\), so \(4r=\frac{1}{m}\), and note that
\[
G(x)=\frac{1-\sqrt{1-4mx}}{2}
\quad\Longrightarrow\quad
G(r)=\frac12.
\]
Recall the inverse map
\[
\Psi_x(v)=\frac{1+v-\sqrt{(1-v)^2+4x}}{2}.
\]
Define the critical iterates
\[
w_0:=G(r)=\frac12,\qquad
w_{k+1}:=\Psi_r(w_k)\quad (k=0,1,\dots,m-2).
\]
Then
\[
\partial_v\Psi_x(v)=\frac12\left(1+\frac{1-v}{\sqrt{(1-v)^2+4x}}\right),
\]
and at \(x=r\) this becomes
\[
\partial_v\Psi_r(v)=\frac12\left(1+\frac{1-v}{\sqrt{(1-v)^2+\frac{1}{m}}}\right).
\]

\begin{proposition}[Limit proportion of $1$'s]\label{prop:p1m-product}
The convergence constant \(\ptop(m)=\lim_{n\to\infty} g_n^{(m-1)}/g_n\) exists and is given by
\[
\boxed{\;
\ptop(m)=\prod_{k=0}^{m-2}\partial_v\Psi_r(w_k)
=\prod_{k=0}^{m-2}
\frac12\left(1+\frac{1-w_k}{\sqrt{(1-w_k)^2+\frac{1}{m}}}\right).
\;}
\]
\end{proposition}

\begin{proof}
Let \(U(x)=G_{m-1}(x)\). By Proposition~\ref{prop:godel-m-psi} we have
\[
U(x)=\Psi_x^{\,m-1}\!\big(G(x)\big),
\qquad\text{hence}\qquad
g_n^{(m-1)}=[x^n]U(x).
\]
The dominant singularity of \(G(x)=\frac{1-\sqrt{1-4mx}}{2}\) is \(r=\frac{1}{4m}\), and with
\(s:=\sqrt{1-x/r}\) (so \(x=r(1-s^2)\)) we have the local expansion
\[
G(x)=\frac12-\frac12\,s+O(s^2)\qquad (x\to r^-).
\]

Define analytic iterates \(w_0(x):=G(x)\) and \(w_{k+1}(x):=\Psi_x(w_k(x))\) for
\(k=0,1,\dots,m-2\); then \(U(x)=w_{m-1}(x)\) and \(w_k=w_k(r)\) are precisely the critical
values defined above.
Since \(\Psi_x(v)\) is analytic in \((x,v)\) near each \((r,w_k)\) and \(x-r=O(s^2)\), a Taylor
expansion shows that if
\[
w_k(x)=w_k-\kappa_k\,s+O(s^2),
\]
then
\[
w_{k+1}(x)=\Psi_x(w_k(x))
=\Psi_r(w_k)-\partial_v\Psi_r(w_k)\,\kappa_k\,s+O(s^2)
=w_{k+1}-\kappa_{k+1}\,s+O(s^2),
\]
where \(\kappa_{k+1}=\partial_v\Psi_r(w_k)\,\kappa_k\).  (Here the \(x\)-variation contributes
only \(O(x-r)=O(s^2)\) and does not affect the linear term in \(s\).)  With \(\kappa_0=\frac12\)
from the expansion of \(G\), iteration gives
\[
\kappa_{m-1}=\frac12\prod_{k=0}^{m-2}\partial_v\Psi_r(w_k).
\]
Consequently,
\[
U(x)=U(r)-\kappa_{m-1}\sqrt{1-\frac{x}{r}}+O\!\left(1-\frac{x}{r}\right)\qquad (x\to r^-).
\]

By the transfer theorem for square--root singularities \cite[Ch.~VI]{FlajoletSedgewick2009},
\[
g_n=[x^n]G(x)\sim \frac{r^{-n}}{4\sqrt{\pi}\,n^{3/2}},
\qquad
g_n^{(m-1)}=[x^n]U(x)\sim \frac{\kappa_{m-1}\,r^{-n}}{2\sqrt{\pi}\,n^{3/2}}.
\]
Taking the ratio gives the existence of the limit and
\[
\ptop(m)=\lim_{n\to\infty}\frac{g_n^{(m-1)}}{g_n}=2\kappa_{m-1}
=\prod_{k=0}^{m-2}\partial_v\Psi_r(w_k).
\]
Substituting \(\partial_v\Psi_r(v)=\frac12\!\left(1+\frac{1-v}{\sqrt{(1-v)^2+\frac1m}}\right)\)
gives the stated explicit product.
\end{proof}

\FloatBarrier
\noindent
\textbf{Numerics.}
Table~\ref{tab:p1m} lists \(\ptop(m)\) for several values of \(m\).  The values decrease with \(m\)
and stabilise as \(m\) becomes large.
\begin{table}[H]
\centering
\tiny
\renewcommand{\arraystretch}{1.05}
\begin{tabular}{r|cccccccc}
\hline
\(m\) & 2 & 3 & 4 & 5 & 6 & 8 & 10 & 20 \\
\hline
\(\ptop(m)\) &
0.7886751346 &
0.7190479224 &
0.6841222107 &
0.6630699362 &
0.6489761394 &
0.6312738082 &
0.6205979134 &
0.5990962874 \\
\hline
\end{tabular}
\caption{The convergence constant \(\ptop(m)=\lim_{n\to\infty} g_n^{(m-1)}/g_n\), i.e.\ the limiting
proportion of outputs equal to \(1\), for Gödel--\(m\) as \(m\) increases.}
\label{tab:p1m}
\end{table}

\begin{table}[H]
\centering
\small
\renewcommand{\arraystretch}{1.05}
\begin{tabular}{r|ccccc}
\hline
\(m\) & 100 & 200 & 1000 & 10000 & \(1/\sqrt{3}\) \\
\hline
\(\ptop(m)\) &
0.5817218321 &
0.5795375421 &
0.5777879655 &
0.5773940443 &
0.5773502692 \\
\hline
\end{tabular}
\caption{Large-\(m\) behaviour of \(\ptop(m)\), showing convergence to \(1/\sqrt{3}\).}
\label{tab:p1m-large}
\end{table}
\FloatBarrier


\noindent
\textbf{Large-\(m\) limit.}
The \(\Psi_r\)-iteration admits a continuum approximation when \(m\) is large (since \(r=1/(4m)\) is small
but the depth is \(m-1\)). This gives the limiting constant \(p_1(\infty)\).

\begin{theorem}[As $m\to\infty$]
\label{thm:p1limit}
Let \(r=\frac1{4m}\) and let \((w_k)_{k\ge0}\) be defined by \(w_{k+1}=\Psi_r(w_k)\) with
\(w_0=\frac12\). Define
\[
\ptop(m):=\prod_{k=0}^{m-2}\partial_v\Psi_r(w_k).
\]
Then
\[
\boxed{\;\;\ptop(m)\longrightarrow \frac{1}{\sqrt{3}}.\;\;}
\]
\end{theorem}

\begin{proof}
Write \(u_k:=1-w_k\). In our situation \(\Psi_r\) admits the closed form
\[
\Psi_r(w)=\frac{1+w-\sqrt{(1-w)^2+\frac1m}}{2},
\qquad\text{so that}\qquad
u_{k+1}=\frac{u_k+\sqrt{u_k^2+\frac1m}}{2},
\]
and
\[
\partial_v\Psi_r(w)
=\frac12\left(1+\frac{u}{\sqrt{u^2+\frac1m}}\right),
\qquad u=1-w.
\]
Since \(u_{k+1}-u_k=\bigl(\sqrt{u_k^2+\frac1m}-u_k\bigr)/2\ge0\), the sequence \((u_k)\) is increasing.
Moreover \(u_0=\frac12\), and if \(u_k\le 1\) then
\[
u_{k+1}\le \frac{u_k+\sqrt{u_k^2+1}}{2}\le \frac{u_k+(u_k+1)}{2}\le 1.
\]
Hence
\[
\frac12\le u_k\le 1\qquad (0\le k\le m).
\]

\medskip
\noindent
\emph{Uniform expansion of \(\log(\partial_v\Psi_r(w_k))\).}
Let \(\varepsilon:=\frac{1}{m u^2}\). For \(u\in[1/2,1]\) we have \(u^2\ge 1/4\), hence
\[
0\le \varepsilon=\frac{1}{m u^2}\le \frac{4}{m},
\]
so \(\varepsilon\to0\) uniformly as \(m\to\infty\). Using
\[
\frac{u}{\sqrt{u^2+\frac1m}}=\frac{1}{\sqrt{1+\varepsilon}}
=1-\frac{\varepsilon}{2}+O(\varepsilon^2),
\]
we obtain, uniformly for \(u\in[1/2,1]\),
\begin{equation}
\label{eq:derivexp_rewrite}
\partial_v\Psi_r(1-u)
=1-\frac{1}{4m u^2}+O\!\left(\frac{1}{m^2}\right).
\end{equation}
Set
\[
x_k:=\frac{1}{4m u_k^2},\qquad \delta_k:=\partial_v\Psi_r(w_k)-\bigl(1-x_k\bigr).
\]
Then \(0\le x_k\le 1/m\) and, by \eqref{eq:derivexp_rewrite}, \(|\delta_k|\le C/m^2\) uniformly for \(k\le m\).
In particular \(x_k-\delta_k=O(1/m)\) uniformly, hence for \(m\) large enough \(\partial_v\Psi_r(w_k)\in(1/2,1)\)
uniformly and we may use \(\log(1-y)=-y+O(y^2)\) uniformly:
\[
\log\bigl(\partial_v\Psi_r(w_k)\bigr)
=\log\bigl(1-(x_k-\delta_k)\bigr)
=-(x_k-\delta_k)+O\bigl((x_k-\delta_k)^2\bigr).
\]
Summing from \(k=0\) to \(m-2\) and using
\[
\sum_{k\le m}x_k^2\le m\cdot (1/m^2)=O(1/m),\qquad
\sum_{k\le m}|\delta_k|=O(1/m),
\]
we obtain
\begin{equation}
\label{eq:logp1_rewrite}
\log \ptop(m)
=\sum_{k=0}^{m-2}\log\bigl(\partial_v\Psi_r(w_k)\bigr)
=-\sum_{k=0}^{m-2}\frac{1}{4m u_k^2}+o(1).
\end{equation}

\medskip
\noindent
\emph{Euler approximation.}
Let \(h:=1/m\) and \(t_k:=kh\). Let \(u(t)\) be the solution of
\[
u'(t)=\frac{1}{4u(t)},\qquad u(0)=\frac12,
\]
so \(u(t)^2=\frac14+\frac{t}{2}\) and \(u(t)\in[1/2,1]\) for \(t\in[0,1]\).
From the exact recurrence,
\[
u_{k+1}-u_k=\frac{\sqrt{u_k^2+h}-u_k}{2}
=\frac{h}{2\bigl(\sqrt{u_k^2+h}+u_k\bigr)}
=\frac{h}{4u_k}+O(h^2),
\]
uniformly for \(u_k\in[1/2,1]\). Thus
\[
u_{k+1}=u_k+h f(u_k)+\rho_k,\qquad f(u)=\frac{1}{4u},\qquad |\rho_k|\le C_1 h^2.
\]
By Taylor's theorem applied to the ODE solution,
\[
u(t_{k+1})=u(t_k)+h f\bigl(u(t_k)\bigr)+\eta_k,\qquad |\eta_k|\le C_2 h^2,
\]
uniformly for \(t_k\in[0,1]\). Setting \(e_k:=u_k-u(t_k)\) and subtracting gives
\[
e_{k+1}=e_k+h\bigl(f(u_k)-f(u(t_k))\bigr)+(\rho_k-\eta_k).
\]
Since \(|f'(u)|=1/(4u^2)\le 1\) on \([1/2,1]\), \(f\) is Lipschitz with constant \(L=1\), hence
\[
|e_{k+1}|\le (1+Lh)|e_k|+C_3 h^2.
\]
Iterating and using \(e_0=0\) yields \(|e_k|\le C_4 h\) uniformly for \(0\le k\le m\), i.e.
\[
\max_{0\le k\le m}|u_k-u(t_k)|=O(h)=O(1/m).
\]

\medskip
\noindent
\emph{Riemann-sum limit.}
Since \(t\mapsto 1/u(t)^2\) is continuous on \([0,1]\) and bounded, and \(u_k=u(t_k)+O(1/m)\) uniformly, we have
\[
\frac{1}{m}\sum_{k=0}^{m-2}\frac{1}{u_k^2}
=\frac{1}{m}\sum_{k=0}^{m-2}\frac{1}{u(t_k)^2}+o(1)
\longrightarrow \int_0^1\frac{dt}{u(t)^2}.
\]
Combining this with \eqref{eq:logp1_rewrite} gives
\[
\log \ptop(m)\longrightarrow -\int_0^1\frac{dt}{4u(t)^2}.
\]
Using \(u(t)^2=\frac14+\frac{t}{2}\),
\[
\int_0^1\frac{dt}{4u(t)^2}
=\int_0^1\frac{dt}{1+2t}
=\frac12\log 3.
\]
Hence \(\log \ptop(m)\to -\frac12\log 3\), i.e.
\[
\ptop(m)\to e^{-(1/2)\log 3}=\frac{1}{\sqrt3},
\]
as claimed.
\end{proof}


\begin{proposition}[Limit proportion of false rows]\label{prop:p0m}
Let \(g_n^{(0)}\) denote the total number of occurrences of the bottom truth value \(v_0=0\)
across all Gödel--\(m\) truth tables for fully bracketed implications on \(n\) variables, and set
\[
\pbot(m):=\lim_{n\to\infty}\frac{g_n^{(0)}}{g_n},\qquad g_n=m^nC_{n-1}.
\]
Then
\[
\boxed{\;
\pbot(m)=\frac12\left(1-\sqrt{\frac{m}{m+4}}\right).
\;}
\]
In particular \(\pbot(m)\sim \frac{1}{m}\) as \(m\to\infty\), so the proportion of false outputs
vanishes in the large--\(m\) limit.
\end{proposition}

\begin{proof}
Let \(H_k(x)=\sum_{j=k}^{m-1}G_j(x)\) be the tail sums, so \(H_0(x)=G(x)\) and \(G_0(x)=H_0(x)-H_1(x)\).
Since \(H_0=\Phi_x(H_1)\) with \(\Phi_x(u)=u+\frac{x}{1-u}\) and \(\Psi_x\) is the inverse branch,
we have \(H_1(x)=\Psi_x(G(x))\). Hence
\[
G_0(x)=G(x)-\Psi_x(G(x)).
\]
At the dominant singularity \(r=\frac{1}{4m}\) we have \(G(r)=\frac12\), and the standard
square--root transfer plus singular--coefficient propagation gives
\[
\pbot(m)=1-\partial_v\Psi_r\!\left(\tfrac12\right).
\]
Using \(\partial_v\Psi_r(v)=\frac12\left(1+\frac{1-v}{\sqrt{(1-v)^2+\frac{1}{m}}}\right)\) gives
\[
\partial_v\Psi_r\!\left(\tfrac12\right)=\frac12\left(1+\frac{\tfrac12}{\sqrt{\tfrac14+\tfrac1m}}\right)
=\frac12\left(1+\sqrt{\frac{m}{m+4}}\right),
\]
which gives the claimed formula.
\end{proof}


\begin{theorem}[Macroscopic cut for nonvanishing limits]\label{thm:cut}
For each \(m\ge2\) and \(k\in\{0,1,\dots,m-1\}\), let
\[
q_k(m):=\lim_{n\to\infty}\frac{[x^n]H_k(x)}{[x^n]G(x)}
=\lim_{n\to\infty}\frac{\#\{\text{outputs }\ge v_k\}}{\#\{\text{all outputs}\}},
\]
so that \(q_{m-1}(m)=\ptop(m)\) and \(p_k(m)=q_k(m)-q_{k+1}(m)\).
Fix \(t\in[0,1]\) and set \(k=\lfloor t(m-1)\rfloor\).  Then, as \(m\to\infty\),
\[
\boxed{\;\;
q_{\lfloor t(m-1)\rfloor}(m)\longrightarrow \frac{1}{\sqrt{\,1+2t\,}}.
\;\;}
\]
Consequently, for any \emph{fixed} level index \(j\) (independent of \(m\)) one has
\(p_j(m)\to0\) as \(m\to\infty\); to obtain a nonzero limit one must place the cut at
a \emph{macroscopic rank} \(k\asymp m\) (equivalently, \(j/(m-1)\to t\)).
\end{theorem}

\begin{proof}
Fix \(t\in[0,1]\) and set \(k=\lfloor t(m-1)\rfloor\). As before, let \(r=\frac{1}{4m}\),
\(w_{j+1}=\Psi_r(w_j)\) with \(w_0=\frac12\), and write \(u_j:=1-w_j\). Then
\[
u_{j+1}=\frac{u_j+\sqrt{u_j^2+\frac1m}}{2},
\qquad
\partial_v\Psi_r(w_j)=\frac12\left(1+\frac{u_j}{\sqrt{u_j^2+\frac1m}}\right).
\]
Since \(u_{j+1}-u_j=\bigl(\sqrt{u_j^2+\frac1m}-u_j\bigr)/2\ge0\), the sequence \((u_j)\) is
increasing, with \(u_0=\frac12\). Moreover, if \(u_j\le 1\) then
\[
u_{j+1}\le \frac{u_j+\sqrt{u_j^2+1}}{2}\le \frac{u_j+(u_j+1)}{2}\le 1,
\]
hence
\[
\frac12\le u_j\le 1\qquad (0\le j\le m).
\]
By the same singular-coefficient propagation as in Proposition~\ref{prop:p1m-product},
applied to the truncated \(\Psi\)-chain defining \(H_k\), one has
\[
q_k(m)=\prod_{j=0}^{k-1}\partial_v\Psi_r(w_j).
\]

\medskip
For \(u\in[1/2,1]\) the uniform expansion \eqref{eq:derivexp_rewrite} gives
\[
\log\bigl(\partial_v\Psi_r(w_j)\bigr)
=-\frac{1}{4m u_j^2}+O\!\left(\frac{1}{m^2}\right),
\]
uniformly for \(0\le j\le m\). Summing from \(j=0\) to \(k-1\) yields
\begin{equation}
\label{eq:logqk_rewrite}
\log q_k(m)
=-\frac{1}{4m}\sum_{j=0}^{k-1}\frac{1}{u_j^2}+o(1),
\end{equation}
since \(k\le m\) implies \(\sum_{j<k}O(m^{-2})=O(k/m^2)=O(1/m)=o(1)\).

\medskip
Let \(h:=1/m\) and \(t_j:=jh\). As above, \((u_j)\) is the Euler scheme for
\(u'(s)=\frac{1}{4u(s)}\), \(u(0)=\frac12\), and the Euler estimate gives
\[
\max_{0\le j\le k}|u_j-u(t_j)|=O(1/m),
\]
where \(u(s)^2=\frac14+\frac{s}{2}\) on \([0,1]\).
Consequently,
\[
\frac{1}{m}\sum_{j=0}^{k-1}\frac{1}{u_j^2}
=\frac{1}{m}\sum_{j=0}^{k-1}\frac{1}{u(t_j)^2}+o(1)
\longrightarrow \int_0^{t}\frac{ds}{u(s)^2},
\]
since \(k/m\to t\). Combining this with \eqref{eq:logqk_rewrite} yields
\[
\log q_k(m)\longrightarrow -\int_0^{t}\frac{ds}{4u(s)^2}
=-\int_0^{t}\frac{ds}{1+2s}
=-\frac12\log(1+2t).
\]
Therefore \(q_{\lfloor t(m-1)\rfloor}(m)\to (1+2t)^{-1/2}\), as claimed.
\end{proof}


\textbf{Remark (all levels \(p_j(m)\)).}
The same singular-coefficient propagation computes \emph{all} convergence constants
\(p_j(m)=\lim_{n\to\infty} g_n^{(j)}/g_n\).  If \(H_k=\sum_{j=k}^{m-1}G_j\) are the tail sums,
then at the critical point \(r=1/(4m)\) one has \(H_k=\Phi_r(H_{k+1})\) with
\(\Phi_x(u)=u+\frac{x}{1-u}\).  Writing \(q_k(m):=\lim_{n\to\infty} [x^n]H_k(x)/[x^n]G(x)\),
one gets the recursion
\[
q_{m-1}(m)=\ptop(m),\qquad
q_k(m)=\big(\partial_u\Phi_r(H_{k+1}(r))\big)\,q_{k+1}(m),
\quad
p_k(m)=q_k(m)-q_{k+1}(m),
\]
where \(\partial_u\Phi_r(u)=1+\frac{r}{(1-u)^2}\).

\paragraph{How the mass spreads across truth values as \(m\) increases.}
For fixed \(m\ge2\) and each \(j\in\{0,1,\dots,m-1\}\), define the \emph{convergence constant}
(limit proportion)
\[
p_j(m)\;:=\;\lim_{n\to\infty}\frac{g_n^{(j)}}{g_n},
\qquad
g_n=m^nC_{n-1}.
\]
Thus \(p_j(m)\) is the asymptotic fraction of outputs equal to \(v_j\) as \(n\to\infty\).
In particular, \(\ptop(m)\) is the limiting proportion of \(1=v_{m-1}\).
These constants satisfy \(\sum_{j=0}^{m-1}p_j(m)=1\), and they determine the asymptotic
coefficients via
\[
g_n^{(j)}\sim \frac{p_j(m)}{4}\cdot \frac{(4m)^n}{\sqrt{\pi}\,n^{3/2}}
\qquad (n\to\infty).
\]

Tables~\ref{tab:pvector-m5} and \ref{tab:pvector-m10} list the full limiting vectors
\((\pbot(m),p_1(m),\dots,p_{m-2}(m), \ptop(m))\) for \(m=5\) and \(m=10\), illustrating how the mass
is distributed among intermediate truth values as \(m\) grows.

\begin{table}[H]
\centering
\small
\renewcommand{\arraystretch}{1.05}
\begin{tabular}{c|ccccc}
\hline
\(j\) (value \(v_j\)) & 0 & 1 & 2 & 3 & 4 \\
\hline
\(p_j(5)\) & 0.127322 & 0.089600 & 0.067242 & 0.052766 & 0.663070 \\
\hline
\end{tabular}
\caption{Limit proportions \(p_j(5)=\lim_{n\to\infty} g_n^{(j)}/g_n\) for Gödel--\(5\).  The
top entry \(p_4(5)\) is the limiting proportion of \(1=v_4\).}
\label{tab:pvector-m5}
\end{table}

\begin{table}[H]
\centering
\tiny
\renewcommand{\arraystretch}{1.05}
\begin{tabular}{c|cccccccccc}
\hline
\(j\) (value \(v_j\)) & 0 & 1 & 2 & 3 & 4 & 5 & 6 & 7 & 8 & 9 \\
\hline
\(p_j(10)\) &
0.077423 & 0.062152 & 0.051275 & 0.043214 & 0.037050 &
0.032213 & 0.028338 & 0.025177 & 0.022560 & 0.620598 \\
\hline
\end{tabular}
\caption{Limit proportions \(p_j(10)=\lim_{n\to\infty} g_n^{(j)}/g_n\) for Gödel--\(10\).  The
top entry \(p_9(10)\) is the limiting proportion of \(1=v_9\).}
\label{tab:pvector-m10}
\end{table}

\newpage

\oneLinePart{ Limit Law as $m\to\infty$}\label{sec:limit-law-shape}

For each $m\ge2$, let
\[
p_k(m):=\lim_{n\to\infty}\frac{g_n^{(k)}}{g_n}\qquad (k=0,1,\dots,m-1)
\]
denote the limiting proportion of output entries equal to $v_k$ in the Gödel--$m$ truth tables.
Equivalently, $(p_k(m))_{k=0}^{m-1}$ is a probability distribution on $\{0,1,\dots,m-1\}$.
Define the tail (survival) sums
\[
q_k(m):=\sum_{j=k}^{m-1}p_j(m)=\Pr(K_m\ge k),
\]
where $K_m$ is a random variable with $\Pr(K_m=k)=p_k(m)$.

\subsection{Macroscopic cuts}

\begin{theorem}[Macroscopic cut for nonvanishing limits]\label{thm:macro-cut}
For each $m\ge 2$ and $k\in\{0,1,\dots,m-1\}$, one has
\[
q_k(m)=\lim_{n\to\infty}\frac{[x^n]H_k(x)}{[x^n]G(x)}
=\lim_{n\to\infty}\frac{\#\{\text{outputs}\ge v_k\}}{\#\{\text{all outputs}\}},
\]
so that $q_{m-1}(m)=p_{m-1}(m)$ and $p_k(m)=q_k(m)-q_{k+1}(m)$.
Fix $t\in[0,1]$ and set $k=\lfloor t(m-1)\rfloor$. Then, as $m\to\infty$,
\[
q_{\lfloor t(m-1)\rfloor}(m)\longrightarrow \frac{1}{\sqrt{1+2t}}.
\]
Consequently, for any fixed level index $j$ (independent of $m$) one has $p_j(m)\to 0$ as
$m\to\infty$; to obtain a nonzero limit one must place the cut at a macroscopic rank $k\simeq m$
(equivalently, $k/(m-1)\to t$).

\end{theorem}

\begin{proof}
Fix $t\in[0,1]$ and set $k:=\lfloor t(m-1)\rfloor$.  Write $w_0=\frac12$ and define the
iterates at the critical point $r=\frac1{4m}$ by
\[
w_{j+1}:=\Psi_r(w_j)\qquad (j\ge0),
\]
and set
\[
u_j:=1-w_j.
\]
Let $h:=1/m$ and $t_j:=jh$.  As in the proof of Theorem~4, the recursion for $(u_j)$ is the
Euler scheme for the ODE
\[
u'(s)=\frac{1}{4u(s)},\qquad u(0)=\frac12,
\]
whose solution satisfies $u(s)^2=\frac14+\frac{s}{2}$ on $[0,1]$.  Moreover, the Euler error
bound gives
\[
\max_{0\le j\le m}\,|u_j-u(t_j)|=O(1/m),
\]
hence also
\[
\max_{0\le j\le m}\,|u_j^2-u(t_j)^2|=O(1/m).
\]
In particular, for $0\le j\le m$ one has
\[
u_j^2=\frac14+\frac{j}{2m}+O(1/m)
\qquad\text{uniformly in }j.
\]

For $u\in[1/2,1]$ we also have the uniform expansion
\[
\partial_v\Psi_r(1-u)=1-\frac{1}{4mu^2}+O\!\left(\frac1{m^2}\right),
\]
hence
\[
\partial_v\Psi_r(w_j)
=1-\frac{1}{4m\,u_j^2}+O\!\left(\frac1{m^2}\right)
\qquad\text{uniformly for }0\le j\le m.
\]
The same singular-coefficient propagation used for the top level yields the product rule
\[
q_k(m)=\prod_{j=0}^{k-1}\partial_v\Psi_r(w_j).
\]
Therefore, using $\log(1-y)=-y+O(y^2)$ uniformly, 
(using $u_j\in[1/2,1]$, so the argument of $\log$ is $1-O(1/m)$ uniformly),
\[
\log q_k(m)
=\sum_{j=0}^{k-1}\log\!\left(1-\frac{1}{4m\,u_j^2}+O\!\left(\frac1{m^2}\right)\right)
=-\frac{1}{4m}\sum_{j=0}^{k-1}\frac{1}{u_j^2}+o(1),
\]
since $k\le m$ implies $\sum_{j<k}O(1/m^2)=O(k/m^2)=O(1/m)=o(1)$.

Now $k=\lfloor t(m-1)\rfloor$ implies $k/m\to t$, and since $u_j=u(t_j)+O(1/m)$ uniformly,
\[
\frac1m\sum_{j=0}^{k-1}\frac{1}{u_j^2}
=\frac1m\sum_{j=0}^{k-1}\frac{1}{u(t_j)^2}+o(1)
\longrightarrow \int_0^t \frac{ds}{u(s)^2}.
\]
Using $u(s)^2=\frac14+\frac{s}{2}$ gives
\[
\log q_k(m)\longrightarrow
-\int_0^t \frac{ds}{4u(s)^2}
=-\int_0^t \frac{ds}{1+2s}
=-\frac12\log(1+2t),
\]
whence $q_k(m)\to (1+2t)^{-1/2}$.  Since $k=\lfloor t(m-1)\rfloor$, this is the claimed limit.
\end{proof}


\subsection{Limit law and limit shape}
Define the rescaled level
\[
T_m := \frac{K_m}{m-1}\in[0,1],
\qquad\text{so that}\qquad
\Pr\!\left(T_m=\frac{k}{m-1}\right)=\Pr(K_m=k)=p_k(m)
\quad (k=0,1,\dots,m-1).
\]
Equivalently, let
\[
\mu_m := \sum_{k=0}^{m-1} p_k(m)\,\delta_{k/(m-1)}
\]
be the law of $T_m$ as a probability measure on $[0,1]$.

\newpage

\begin{theorem}[Limit law / weak limit of the scaled output level]\label{thm:limit-law-shape}
Assume the macroscopic cut limit of Theorem~6 holds. For $t\in[0,1)$ and let
$k_m(t)=\lceil t(m-1)\rceil$. Then
\begin{equation}\label{eq:survival-Tm}
\Pr(T_m\ge t)=\sum_{j=k_m(t)}^{m-1} p_j(m)=q_{k_m(t)}(m)\longrightarrow \frac1{\sqrt{1+2t}}.
\end{equation}
Assume also that the top mass converges:
\begin{equation}\label{eq:top-mass}
\Pr(T_m=1)=p_{m-1}(m)\longrightarrow \frac1{\sqrt3}.
\end{equation}
Then $T_m\Rightarrow T$ (equivalently $\mu_m\Rightarrow\mu$), where
\begin{equation}\label{eq:limit-law}
\mu(dx)=\frac{1}{(1+2x)^{3/2}}\mathbf 1_{[0,1)}(x)\,dx+\frac1{\sqrt3}\,\delta_1(dx).
\end{equation}
Equivalently, $\Pr(T\ge t)=1/\sqrt{1+2t}$ for $0\le t<1$, and $\Pr(T=1)=1/\sqrt3$.
\end{theorem}

\begin{proof}
Fix $t\in[0,1)$. By~\eqref{eq:survival-Tm} we have
\[
\Pr(T_m\ge t)\longrightarrow (1+2t)^{-1/2}.
\]
Hence the distribution functions
\[
F_m(t):=\Pr(T_m\le t)=1-\Pr(T_m>t)
\]
satisfy $F_m(t)\to F(t)$ for every $t\in[0,1)$, where
\[
F(t):=1-\frac{1}{\sqrt{1+2t}}.
\]
Since $F$ is continuous on $[0,1)$, this determines the weak limit on $[0,1)$.

Assume in addition that the top mass converges, i.e.~\eqref{eq:top-mass} holds:
\[
\Pr(T_m=1)=p_{m-1}(m)\longrightarrow \frac{1}{\sqrt3}.
\]
Thus the limit law has an atom of size $1/\sqrt3$ at $t=1$.

To identify the absolutely continuous part on $[0,1)$, note that $F$ is differentiable on $[0,1)$ and
\[
f(t):=F'(t)=\frac{1}{(1+2t)^{3/2}}.
\]
Moreover,
\[
\int_0^1 f(t)\,dt=\int_0^1 \frac{dt}{(1+2t)^{3/2}}
=\Bigl[-(1+2t)^{-1/2}\Bigr]_{0}^{1}
=1-\frac{1}{\sqrt3},
\]
so the remaining mass is $1/\sqrt3$ and sits at $t=1$. Therefore $T_m\Rightarrow T$
(equivalently $\mu_m\Rightarrow\mu$), where
\[
\mu(dt)=\frac{1}{(1+2t)^{3/2}}\mathbf 1_{[0,1)}(t)\,dt+\frac{1}{\sqrt3}\,\delta_1(dt).
\]
\end{proof}

\begin{corollary}[A derived statistic]
With $T$ as in Theorem~7, one has $\mathbb{E}[T]=\sqrt{3}-1$.
\end{corollary}

\begin{proof}
By Theorem~7 (limit law \eqref{eq:limit-law}), $T$ has density
\[
f(t)=\frac{1}{(1+2t)^{3/2}}\mathbf 1_{[0,1)}(t)
\]
and an atom of mass $1/\sqrt3$ at $t=1$. Hence
\[
\mathbb{E}[T]=\int_0^1 \frac{t}{(1+2t)^{3/2}}\,dt+\frac{1}{\sqrt3}.
\]
With the substitution $u=1+2t$ (so $t=(u-1)/2$ and $dt=du/2$), we get
\[
\int_0^1 \frac{t}{(1+2t)^{3/2}}\,dt
=\frac14\int_1^3 (u^{-1/2}-u^{-3/2})\,du
=\frac12\Bigl(\sqrt{u}+u^{-1/2}\Bigr)\Big|_{1}^{3}
=\frac{2}{\sqrt3}-1.
\]
Therefore
\[
\mathbb{E}[T]=\left(\frac{2}{\sqrt3}-1\right)+\frac{1}{\sqrt3}
=\frac{3}{\sqrt3}-1
=\sqrt3-1.
\]
\end{proof}


\begin{remark}[Equivalent weak convergence formulation]
The convergence $T_m \Rightarrow T$ in Theorem~7 is equivalent to weak convergence of
measures $\mu_m \Rightarrow \mu$, i.e.
\[
\int_{[0,1]} \varphi\, d\mu_m
=\sum_{k=0}^{m-1} p_k(m)\,\varphi\!\left(\frac{k}{m-1}\right)
\longrightarrow
\int_{[0,1]} \varphi\, d\mu,
\]
for every bounded continuous function $\varphi:[0,1]\to\mathbb{R}$, where
\[
\mu(dt)=\frac{1}{(1+2t)^{3/2}}\mathbf 1_{[0,1)}(t)\,dt+\frac{1}{\sqrt3}\,\delta_1(dt).
\]
\end{remark}


\section{  $\;$ A black-box singularity propagation principle}\label{app:blackbox}


\begin{theorem}[Propagation of a square--root singularity under analytic iteration]\label{thm:blackbox-sqrt}
Let $r>0$. Suppose $H_0(x)$ is analytic in a $\Delta$--domain at~$r$ and admits the local expansion
\begin{equation}\label{eq:H0-sqrt}
H_0(x)=a_0-c_0\sqrt{1-\frac{x}{r}}+O\!\left(1-\frac{x}{r}\right)\qquad (x\to r),
\end{equation}
with $c_0\neq 0$.
Let $\Psi(x,v)$ be analytic in a neighbourhood of the compact set
\[
\mathcal{K}:=\{(r,a_k):0\le k\le K\}\subset\mathbb{C}^2,
\]
where $a_0$ is as in \eqref{eq:H0-sqrt} and the real numbers $a_k$ are defined recursively by
\begin{equation}\label{eq:ak-rec}
a_{k+1}=\Psi(r,a_k)\qquad (0\le k\le K-1).
\end{equation}
Assume further that
\begin{equation}\label{eq:dPsi-nonzero}
\partial_v\Psi(r,a_k)\neq 0\qquad (0\le k\le K-1).
\end{equation}
Define $H_k(x)$ recursively by
\begin{equation}\label{eq:Hk-rec}
H_{k+1}(x):=\Psi\!\bigl(x,H_k(x)\bigr)\qquad (0\le k\le K-1).
\end{equation}
Then for each $k\in\{0,1,\dots,K\}$, the function $H_k(x)$ is analytic in a $\Delta$--domain at~$r$
and has a square--root expansion
\begin{equation}\label{eq:Hk-sqrt}
H_k(x)=a_k-c_k\sqrt{1-\frac{x}{r}}+O\!\left(1-\frac{x}{r}\right)\qquad (x\to r),
\end{equation}
where the coefficients satisfy
\begin{equation}\label{eq:ck-rec}
c_{k+1}=c_k\,\partial_v\Psi(r,a_k)\qquad (0\le k\le K-1).
\end{equation}
In particular,
\begin{equation}\label{eq:ck-product}
c_k=c_0\prod_{j=0}^{k-1}\partial_v\Psi(r,a_j)\neq 0\qquad (1\le k\le K).
\end{equation}
Consequently, the coefficient asymptotics are universal:
\begin{equation}\label{eq:coeff-asympt}
[x^n]H_k(x)\sim \frac{c_k}{2\sqrt{\pi}}\;r^{-n}n^{-3/2}\qquad (n\to\infty).
\end{equation}
\end{theorem}

\begin{proof}
Write $t:=1-x/r$, so $t\to 0$ as $x\to r$ (in the $\Delta$--domain). Then \eqref{eq:H0-sqrt} reads
\[
H_0(x)=a_0-c_0\sqrt{t}+O(t),
\]
where $\sqrt{t}=\sqrt{1-x/r}$ denotes the branch analytic in the $\Delta$--domain at $r$.
Fix $k\in\{0,\dots,K-1\}$ and assume inductively that $H_k$ is $\Delta$--analytic at $r$ and has
\eqref{eq:Hk-sqrt} with some $a_k$ and $c_k\neq 0$. We show the same for $H_{k+1}$.

Since $\Psi$ is analytic near $(r,a_k)$, we have the first-order expansion
\[
\Psi(x,v)=\Psi(r,a_k)+\partial_x\Psi(r,a_k)(x-r)+\partial_v\Psi(r,a_k)(v-a_k)+R(x,v),
\]
where
\[
R(x,v)=O\!\left((x-r)^2+(v-a_k)^2+|x-r|\,|v-a_k|\right)
\]
uniformly for $(x,v)$ near $(r,a_k)$.
Substitute $v=H_k(x)$ and use $x-r=-rt$ and $H_k(x)-a_k=-c_k\sqrt{t}+O(t)$ to obtain
\[
H_{k+1}(x)=\underbrace{\Psi(r,a_k)}_{=:a_{k+1}}
-\underbrace{c_k\,\partial_v\Psi(r,a_k)}_{=:c_{k+1}}\,\sqrt{t}+O(t),
\]
which is \eqref{eq:Hk-sqrt} for $k+1$ with $a_{k+1}$ as in \eqref{eq:ak-rec} and $c_{k+1}$ as in
\eqref{eq:ck-rec}. Condition \eqref{eq:dPsi-nonzero} implies $c_{k+1}\neq 0$, so the square--root term
persists. The $\Delta$--analyticity of $H_{k+1}(x)=\Psi(x,H_k(x))$ follows from analyticity of $\Psi$
near $\mathcal{K}$ and the $\Delta$--analyticity of $H_k$ at $r$ (possibly after shrinking the
$\Delta$--domain).

Iterating \eqref{eq:ck-rec} yields \eqref{eq:ck-product}. Finally, the transfer theorem for square--root
singularities applied to \eqref{eq:Hk-sqrt} gives \eqref{eq:coeff-asympt}.
\end{proof}




\oneLinePart{The Sequence Factory}

One advantage of the generating function approach is that it naturally produces, from the same underlying logical model, a large family of structured counting sequences, (we have shown this in Kleene case too, see ~\cite{Yildiz2020}).  Beyond the base output counts (how often each truth value occurs), we can \emph{refine} the enumeration by recording additional local information at the root split of a bracketing, namely the ordered pair of truth values carried by the left and right subformulae.  This refinement decomposes the global truth--table counts into many ``pair--count'' sequences whose ordinary generating functions are simple products of the level OGFs.  In this sense, Gödel--$m$ implication acts as a \emph{sequence factory}: a uniform mechanism that generates an entire collection of natural integer sequences, related by clear algebraic identities and by universal square--root asymptotics.

\begin{table}[h]
\centering
\scriptsize
\setlength{\tabcolsep}{3pt}
\renewcommand{\arraystretch}{1.05}
\caption{Gödel--4 sequences (first ten terms). Base sequences and refined pair--sequences.}
\label{tab:godel4-all}
\resizebox{\textwidth}{!}{%
\begin{tabular}{l|rrrrrrrrrr}
\hline
 & 1 & 2 & 3 & 4 & 5 & 6 & 7 & 8 & 9 & 10\\
\hline
\multicolumn{11}{l}{\textbf{Base sequences}}\\
\hline
$t_n$ & 1 & 10 & 80 & 825 & 9355 & 113237 & 1431976 & 18696855 & 250122284 & 3410617188\\
$b_n$ & 1 & 1 & 11 & 101 & 1116 & 13186 & 164093 & 2116676 & 28052479 & 379673121\\
$a_n$ & 1 & 2 & 15 & 143 & 1559 & 18379 & 228175 & 2938786 & 38902987 & 526061137\\
$f_n$ & 1 & 3 & 22 & 211 & 2306 & 27230 & 338444 & 4362627 & 57788170 & 781825066\\
$g_n$ & 4 & 16 & 128 & 1280 & 14336 & 172032 & 2162688 & 28114944 & 374865920 & 5098176512\\
\hline
\multicolumn{11}{l}{\textbf{Refined false-row pair--sequences}}\\
\hline
$N^{a,0}_n$ & 0 & 1 & 5 & 43 & 443 & 5046 & 61209 & 775277 & 10134321 & 135696192\\
$N^{b,0}_n$ & 0 & 1 & 4 & 36 & 367 & 4178 & 50613 & 640554 & 8368030 & 111992206\\
$N^{1,0}_n$ & 0 & 1 & 13 & 132 & 1496 & 18006 & 226622 & 2946796 & 39285819 & 534136668\\
\hline
\multicolumn{11}{l}{\textbf{Refined $a$- and $b$-row pair--sequences}}\\
\hline
$N^{b,a}_n$ & 0 & 1 & 3 & 28 & 281 & 3185 & 38444 & 485351 & 6328829 & 84580645\\
$N^{1,a}_n$ & 0 & 1 & 12 & 115 & 1278 & 15194 & 189731 & 2453435 & 32574158 & 441480492\\
$N^{1,b}_n$ & 0 & 1 & 11 & 101 & 1116 & 13186 & 164093 & 2116676 & 28052479 & 379673121\\
\hline
\multicolumn{11}{l}{\textbf{Refined true-row pair--sequences ($u\le v$)}}\\
\hline
$N^{0,0}_n$ & 0 & 1 & 6 & 53 & 554 & 6362 & 77580 & 986253 & 12927170 & 173452334\\
$N^{0,a}_n$ & 0 & 1 & 5 & 43 & 443 & 5046 & 61209 & 775277 & 10134321 & 135696192\\
$N^{0,b}_n$ & 0 & 1 & 4 & 36 & 367 & 4178 & 50613 & 640554 & 8368030 & 111992206\\
$N^{0,1}_n$ & 0 & 1 & 13 & 132 & 1496 & 18006 & 226622 & 2946796 & 39285819 & 534136668\\
$N^{a,a}_n$ & 0 & 1 & 4 & 34 & 346 & 3915 & 47284 & 597085 & 7787516 & 104093243\\
$N^{a,b}_n$ & 0 & 1 & 3 & 28 & 281 & 3185 & 38444 & 485351 & 6328829 & 84580645\\
$N^{a,1}_n$ & 0 & 1 & 12 & 115 & 1278 & 15194 & 189731 & 2453435 & 32574158 & 441480492\\
$N^{b,b}_n$ & 0 & 1 & 2 & 23 & 224 & 2555 & 30826 & 389311 & 5077062 & 67857384\\
$N^{b,1}_n$ & 0 & 1 & 11 & 101 & 1116 & 13186 & 164093 & 2116676 & 28052479 & 379673121\\
$N^{1,1}_n$ & 0 & 1 & 20 & 260 & 3250 & 41610 & 545574 & 7306117 & 99586900 & 1377654903\\
\hline
\end{tabular}}
\vspace{2pt}
\par\footnotesize\noindent
Note that $N^{u,v}_1=0$ for all pairs $(u,v)$ since there is no binary split at size $n=1$.
\end{table}



\newpage



\section{  $\;$ Refined pair-count sequences in the G\"odel--4 case}\label{app:godel4-refined-sequences}

As in the Kleene case~\cite{Yildiz2020}, one may refine the basic output sequences by recording,
at the \emph{root split} $\varphi \Rightarrow_G \psi$, the ordered pair of truth values
$\bigl(\nu(\varphi),\nu(\psi)\bigr)$.
This refinement is particularly transparent in the G\"odel--$4$ case, where the truth values are
\[
0<a<b<1,
\]
and the G\"odel implication is
\[
u\Rightarrow_G w=
\begin{cases}
1, & u\le w,\\
w, & u>w.
\end{cases}
\]

Let $F(x),A(x),B(x),T(x)$ be the OGFs counting rows whose \emph{output} value is $0,a,b,1$, respectively,
so
\[
F(x)=\sum_{n\ge1} f_n x^n,\quad
A(x)=\sum_{n\ge1} a_n x^n,\quad
B(x)=\sum_{n\ge1} b_n x^n,\quad
T(x)=\sum_{n\ge1} t_n x^n,
\]
with $f_1=a_1=b_1=t_1=1$ (a single variable can take each truth value once).

\medskip
\noindent\textbf{Pair-count refinement.}
For $u,w\in\{0,a,b,1\}$ and $n\ge2$, define $N^{u,w}_n$ to be the total number of truth-table rows
\emph{across all fully bracketed implications on $n$ variables and all valuations} whose root split
$\varphi\Rightarrow_G\psi$ satisfies
\[
\nu(\varphi)=u,\qquad \nu(\psi)=w.
\]
(For convenience set $N^{u,w}_1:=0$.)  Let $U(x)\in\{F(x),A(x),B(x),T(x)\}$ be the OGF corresponding to $u$,
and similarly let $W(x)$ correspond to $w$. Then for $n\ge2$,
\[
N^{u,w}_n=\sum_{i=1}^{n-1} U_i\,W_{n-i},
\qquad\text{so}\qquad
N^{u,w}(x):=\sum_{n\ge2}N^{u,w}_n x^n = U(x)\,W(x).
\]
Thus the pair-count OGFs are \emph{simple products} of the four base OGFs.

\medskip
\noindent\textbf{Recovering the output OGFs by summing pair classes.}
Since the output of $u\Rightarrow_G w$ depends only on the comparison of $u$ and $w$, the pair-count OGFs
immediately decompose $F,A,B,T$.

\medskip
\noindent\emph{False output $0$.}
This occurs iff $w=0$ and $u>0$, i.e.\ $(u,w)\in\{(a,0),(b,0),(1,0)\}$. Hence, for $n\ge2$,
\[
f_n = N^{a,0}_n+N^{b,0}_n+N^{1,0}_n,
\]
and at the OGF level
\[
F(x)-x = A(x)F(x)+B(x)F(x)+T(x)F(x)=\bigl(A(x)+B(x)+T(x)\bigr)F(x).
\]

\medskip
\noindent\emph{Output $a$.}
This occurs iff $w=a$ and $u>a$, i.e.\ $(u,w)\in\{(b,a),(1,a)\}$. Thus
\[
A(x)-x = B(x)A(x)+T(x)A(x)=\bigl(B(x)+T(x)\bigr)A(x).
\]

\medskip
\noindent\emph{Output $b$.}
This occurs iff $w=b$ and $u>b$, i.e.\ only $(u,w)=(1,b)$. Thus
\[
B(x)-x = T(x)B(x).
\]

\medskip
\noindent\emph{True output $1$.}
This occurs iff $u\le w$. There are $10$ admissible pairs:
\[
(0,0),(0,a),(0,b),(0,1),\ (a,a),(a,b),(a,1),\ (b,b),(b,1),\ (1,1),
\]
hence
\[
T(x)-x=\!\!\sum_{(u,w):\,u\le w}\! U(x)W(x)
=F^2+FA+FB+FT+A^2+AB+AT+B^2+BT+T^2.
\]

\medskip
In summary, see ~\ref{tab:godel4-all}, the $16$ pair-count sequences $\{N_n^{u,w}\}$ (or $16$ pair-count OGFs $N^{u,w}(x)$) form a
finest root-split decomposition; the four output sequences $(f_n,a_n,b_n,t_n)$ are recovered by summing the
appropriate subclasses dictated by the G\"odel implication rule.


\section{ $\;$ Refined pair-count sequences in the G\"odel--$m$ case}\label{app:godelm-pair-sequences}

Let the truth values be totally ordered
\[
v_0< v_1<\cdots< v_{m-1}=1,
\]
and let $G_k(x)=\sum_{n\ge1} g^{(k)}_n x^n$ denote the OGF counting rows whose \emph{output} value is exactly $v_k$.
Let
\[
G(x)=\sum_{k=0}^{m-1} G_k(x)=\frac{1-\sqrt{1-4mx}}{2},
\]
so the dominant singularity is $r=\frac1{4m}$.

\begin{theorem}[Number of refined pair--sequences]\label{thm:pair-sequences}
For each ordered pair $(i,j)\in\{0,\dots,m-1\}^2$ and each $n\ge 1$, let $N_n^{i,j}$ be the total number of
truth--table rows \emph{across all fully bracketed implications on $n$ variables and all valuations}
for which the root split $\varphi \Rightarrow_G \psi$ satisfies
\[
\nu(\varphi)=v_i \qquad\text{and}\qquad \nu(\psi)=v_j.
\]
(Equivalently, set $N_1^{i,j}:=0$ and count root splits for $n\ge2$.)
Then:
\begin{enumerate}
\item[(i)] There are exactly $m^2$ refined pair--sequences $\{N_n^{i,j}\}_{n\ge 1}$.
\item[(ii)] Their OGFs satisfy, for all $(i,j)$,
\[
N^{i,j}(x):=\sum_{n\ge 1} N_n^{i,j}x^n \;=\; G_i(x)\,G_j(x),
\qquad\text{and in particular } N_1^{i,j}=0.
\]
\item[(iii)] Under G\"odel implication, $v_i \Rightarrow_G v_j = v_{m-1}=1$ if $i\le j$, and
$v_i \Rightarrow_G v_j = v_j$ if $i>j$. Hence $\frac{m(m+1)}{2}$ of the pair--sequences (those with $i\le j$)
contribute to the top value $v_{m-1}$, while for each $j\le m-2$ exactly $m-1-j$ pair--sequences
(those with fixed right index $j$ and $i>j$) contribute to the output value $v_j$.
\end{enumerate}
\end{theorem}

\begin{proof}
For $n\ge2$, every bracketing of size $n$ splits uniquely at the root into a left subformula of size $t$
and a right subformula of size $n-t$, where $1\le t\le n-1$. Counting rows with $\nu(\varphi)=v_i$ and
$\nu(\psi)=v_j$ gives the convolution
\[
N_n^{i,j}=\sum_{t=1}^{n-1} g_t^{(i)}\,g_{n-t}^{(j)}.
\]
Summing over $n\ge1$ yields $N^{i,j}(x)=G_i(x)G_j(x)$. The count in (i) is immediate since there are $m$ choices
for $i$ and $m$ choices for $j$. Finally, (iii) is exactly the defining rule of G\"odel implication.
\end{proof}

\begin{remark}[Refined pairs versus derived output sequences]\label{rem:pairs-vs-outputs}
The pair--sequences $\{N_n^{i,j}\}$ form the finest root-split decomposition. All output-count sequences
$g_n^{(j)}$ are obtained by summing appropriate $N_n^{i,j}$ according to the G\"odel implication rule:
for $j\le m-2$,
\[
g_n^{(j)}=\sum_{i=j+1}^{m-1} N_n^{i,j},
\]
while the top level is
\[
g_n^{(m-1)}=\sum_{0\le i\le j\le m-1} N_n^{i,j}.
\]
Thus the output sequences are \emph{not} additional independent families beyond $\{N_n^{i,j}\}$.
\end{remark}


\begin{theorem}[Uniform square-root asymptotics for all pair--sequences]\label{thm:pair-uniform-asympt}
For each $k\in\{0,\dots,m-1\}$ there exist constants $a_k\in\mathbb{R}$ and $b_k\neq0$ such
that, as $x\to r^-$,
\[
G_k(x)=a_k-b_k\sqrt{1-\frac{x}{r}}+O\!\left(1-\frac{x}{r}\right).
\]
Consequently,
\[
[x^n]\,G_k(x)\sim \frac{b_k}{2\sqrt{\pi}}\; r^{-n}\,n^{-3/2}
=\frac{b_k}{2\sqrt{\pi}}\,(4m)^n\,n^{-3/2}.
\]
Moreover, for every ordered pair $(i,j)$ one has
\[
N^{i,j}(x)=A_{i,j}-B_{i,j}\sqrt{1-\frac{x}{r}}+O\!\left(1-\frac{x}{r}\right),
\qquad
B_{i,j}=a_i b_j+a_j b_i,
\]
and hence
\[
N^{i,j}_n \sim \frac{B_{i,j}}{2\sqrt{\pi}}\; r^{-n}\,n^{-3/2}
      \;=\; \frac{B_{i,j}}{2\sqrt{\pi}}\,(4m)^n\,n^{-3/2}.
\]
In particular, all $m^2$ refined pair--sequences share the same dominant singularity $r$,
the same exponential growth $(4m)^n$, and the same polynomial decay $n^{-3/2}$; only the
constants depend on $(i,j)$.
\end{theorem}

\begin{proof}
Let $H_\ell(x):=\sum_{k=\ell}^{m-1}G_k(x)$ be the tail sums. Then $H_0(x)=G(x)$ and
\[
H_{\ell+1}(x)=\Psi_x\bigl(H_\ell(x)\bigr)\qquad(\ell=0,1,\dots,m-2),
\]
where $\Psi_x$ is the inverse branch of $\Phi_x(u)=u+\frac{x}{1-u}$, namely
\[
\Psi_x(v)=\frac{1+v-\sqrt{(1-v)^2+4x}}{2}.
\]
Write $\Delta(x)=1-\frac{x}{r}$. Since
\[
H_0(x)=G(x)=\frac12-\frac12\sqrt{\Delta(x)},
\]
$H_0$ has a square-root expansion at $r$. The analytic-iteration theorem
(Theorem~\ref{thm:blackbox-sqrt}) applied to $H_{\ell+1}=\Psi_x(H_\ell)$ propagates the
square-root form to all $H_\ell$:
\[
H_\ell(x)=h_\ell-c_\ell\sqrt{\Delta(x)}+O(\Delta(x)),
\qquad c_\ell\neq0.
\]
Then $G_k(x)=H_k(x)-H_{k+1}(x)$ for $k\le m-2$ and $G_{m-1}(x)=H_{m-1}(x)$, so each $G_k$
inherits a square-root expansion
\[
G_k(x)=a_k-b_k\sqrt{\Delta(x)}+O(\Delta(x)),
\]
with $b_k\neq0$. Multiplying $G_i(x)G_j(x)$ gives the stated expansion for $N^{i,j}(x)$ with
$B_{i,j}=a_i b_j+a_j b_i$. The coefficient asymptotics follow from the standard transfer
theorem for square-root singularities~\cite[Theorem~VI.1]{FlajoletSedgewick2009}.
\end{proof}


\begin{proposition}[Symmetry and an upper bound on distinct products]\label{prop:pair-seq-symmetry}
Fix $m\ge2$ and define $N^{i,j}(x)=G_i(x)G_j(x)$ as above.
\begin{enumerate}
\item[(i)] (Swap symmetry) For all $i,j$ one has $N^{i,j}(x)=N^{j,i}(x)$.
\item[(ii)] (Upper bound) Among the $m^2$ ordered pairs $(i,j)$ there are at most
\[
\binom{m+1}{2}=\frac{m(m+1)}{2}
\]
distinct OGFs $N^{i,j}(x)$ (equivalently, at most that many distinct integer sequences $N^{i,j}_n$),
since swap symmetry identifies $(i,j)$ with $(j,i)$.
\end{enumerate}
\end{proposition}

\begin{proof}
(i) is commutativity of multiplication. (ii) follows because unordered pairs with repetition are counted by
$\binom{m+1}{2}$.
\end{proof}

\begin{remark}[No cross--$m$ coincidences]\label{rem:no-cross-m}
Sequences coming from different values of $m$ cannot coincide: by Theorem~\ref{thm:pair-uniform-asympt},
every pair-sequence has exponential growth $(4m)^n$. Hence a sequence from Gödel--$m$ cannot equal a sequence
from Gödel--$m'$ if $m\neq m'$.
\end{remark}

\begin{remark}[A small ``sequence factory'']\label{rem:sequence-factory}
For each $m$ and each ordered pair $(v_i,v_j)$, the Gödel--$m$ implication scheme produces a natural
integer sequence $N_n^{i,j}$ whose OGF is the simple product $G_i(x)G_j(x)$ and whose coefficients satisfy
a universal square-root asymptotic law
\[
N^{i,j}_n \sim C_{i,j}\,(4m)^n\,n^{-3/2}.
\]
\end{remark}





\newpage
\appendix

\section{Supplementary material}

\begin{table}[h]
\centering
\[
\begin{array}{c|c|c}
p_1 & p_2 & p_1 \Rightarrow_G p_2\\
\hline
0 & 0 & 1\\
0 & a & 1\\
0 & b & 1\\
0 & 1 & 1\\
a & 0 & 0\\
a & a & 1\\
a & b & 1\\
a & 1 & 1\\
b & 0 & 0\\
b & a & a\\
b & b & 1\\
b & 1 & 1\\
1 & 0 & 0\\
1 & a & a\\
1 & b & b\\
1 & 1 & 1
\end{array}
\]
\caption{Explicit truth table for \(p_1 \Rightarrow_G p_2\) on \(V_4\).}
\label{tab:godel-n2}
\end{table}

From Table~8 we obtain the distribution of truth values in the output column of the truth table.
There are three occurrences of the value zero, two occurrences of the value \(a\), one occurrence
of the value \(b\), and ten occurrences of the value one. Consequently, the total number of
entries is sixteen.

\newpage
\begin{table}[ht]
\centering
\caption{Truth tables for the two bracketings on three variables in Gödel four-valued logic.}
\label{tab:godel4-n3-sidebyside}

\begin{minipage}{0.3\textwidth}
\centering
\resizebox{0.60\linewidth}{!}{%
\begin{tabular}{cccc}
\hline
$p_1$ & $p_2$ & $p_3$ & ($p_1\Rightarrow_G p_2) \Rightarrow_G p_3$\\
\hline
0 & 0 & 0 & 0 \\
0 & 0 & a & a \\
0 & 0 & b & b \\
0 & 0 & 1 & 1 \\
0 & a & 0 & 0 \\
0 & a & a & a \\
0 & a & b & b \\
0 & a & 1 & 1 \\
0 & b & 0 & 0 \\
0 & b & a & a \\
0 & b & b & b \\
0 & b & 1 & 1 \\
0 & 1 & 0 & 0 \\
0 & 1 & a & a \\
0 & 1 & b & b \\
0 & 1 & 1 & 1 \\
a & 0 & 0 & 1 \\
a & 0 & a & 1 \\
a & 0 & b & 1 \\
a & 0 & 1 & 1 \\
a & a & 0 & 0 \\
a & a & a & a \\
a & a & b & b \\
a & a & 1 & 1 \\
a & b & 0 & 0 \\
a & b & a & a \\
a & b & b & b \\
a & b & 1 & 1 \\
a & 1 & 0 & 0 \\
a & 1 & a & a \\
a & 1 & b & b \\
a & 1 & 1 & 1 \\
b & 0 & 0 & 1 \\
b & 0 & a & 1 \\
b & 0 & b & 1 \\
b & 0 & 1 & 1 \\
b & a & 0 & 0 \\
b & a & a & 1 \\
b & a & b & 1 \\
b & a & 1 & 1 \\
b & b & 0 & 0 \\
b & b & a & a \\
b & b & b & b \\
b & b & 1 & 1 \\
b & 1 & 0 & 0 \\
b & 1 & a & a \\
b & 1 & b & b \\
b & 1 & 1 & 1 \\
1 & 0 & 0 & 1 \\
1 & 0 & a & 1 \\
1 & 0 & b & 1 \\
1 & 0 & 1 & 1 \\
1 & a & 0 & 0 \\
1 & a & a & 1 \\
1 & a & b & 1 \\
1 & a & 1 & 1 \\
1 & b & 0 & 0 \\
1 & b & a & a \\
1 & b & b & 1 \\
1 & b & 1 & 1 \\
1 & 1 & 0 & 0 \\
1 & 1 & a & a \\
1 & 1 & b & b \\
1 & 1 & 1 & 1 \\
\hline
\end{tabular}%
}
\end{minipage}
\hfill
\begin{minipage}{0.3\textwidth}
\centering
\resizebox{0.60\linewidth}{!}{%
\begin{tabular}{cccc}
\hline
$p_1$ & $p_2$ & $p_3$ & $p_1\Rightarrow_G (p_2\Rightarrow_G p_3)$\\
\hline
0 & 0 & 0 & 1 \\
0 & 0 & a & 1 \\
0 & 0 & b & 1 \\
0 & 0 & 1 & 1 \\
0 & a & 0 & 1 \\
0 & a & a & 1 \\
0 & a & b & 1 \\
0 & a & 1 & 1 \\
0 & b & 0 & 1 \\
0 & b & a & 1 \\
0 & b & b & 1 \\
0 & b & 1 & 1 \\
0 & 1 & 0 & 1 \\
0 & 1 & a & 1 \\
0 & 1 & b & 1 \\
0 & 1 & 1 & 1 \\
a & 0 & 0 & 1 \\
a & 0 & a & 1 \\
a & 0 & b & 1 \\
a & 0 & 1 & 1 \\
a & a & 0 & 0 \\
a & a & a & 1 \\
a & a & b & 1 \\
a & a & 1 & 1 \\
a & b & 0 & 0 \\
a & b & a & 1 \\
a & b & b & 1 \\
a & b & 1 & 1 \\
a & 1 & 0 & 0 \\
a & 1 & a & 1 \\
a & 1 & b & 1 \\
a & 1 & 1 & 1 \\
b & 0 & 0 & 1 \\
b & 0 & a & 1 \\
b & 0 & b & 1 \\
b & 0 & 1 & 1 \\
b & a & 0 & 0 \\
b & a & a & 1 \\
b & a & b & 1 \\
b & a & 1 & 1 \\
b & b & 0 & 0 \\
b & b & a & a \\
b & b & b & 1 \\
b & b & 1 & 1 \\
b & 1 & 0 & 0 \\
b & 1 & a & a \\
b & 1 & b & 1 \\
b & 1 & 1 & 1 \\
1 & 0 & 0 & 1 \\
1 & 0 & a & 1 \\
1 & 0 & b & 1 \\
1 & 0 & 1 & 1 \\
1 & a & 0 & 0 \\
1 & a & a & 1 \\
1 & a & b & 1 \\
1 & a & 1 & 1 \\
1 & b & 0 & 0 \\
1 & b & a & a \\
1 & b & b & 1 \\
1 & b & 1 & 1 \\
1 & 1 & 0 & 0 \\
1 & 1 & a & a \\
1 & 1 & b & b \\
1 & 1 & 1 & 1 \\
\hline
\end{tabular}%
}
\end{minipage}

\end{table}




\newpage

$\;$\\

\begin{thebibliography}{99}

\bibitem{CameronYildiz2010}
P.~J.~Cameron and V.~Yildiz,
\newblock \emph{Counting false entries in truth tables of bracketed formulae connected by implication},
\newblock arXiv:1106.4443 (2010).

\bibitem{Yildiz2011b}
V.~Yildiz,
\newblock \emph{Counting false entries in truth tables of bracketed formulae connected by $m$-implication},
\newblock arXiv:1203.4645 (2011).

\bibitem{Yildiz2012}
V.~Yildiz,
\newblock \emph{General Combinatorical Structure of Truth Tables Connected by Implication},
\newblock arXiv:1205.5595 (2012).

\bibitem{Yildiz2020}
V.~Yildiz,
\newblock \emph{Counting entries in Kleene truth table},
\newblock arXiv:2010.10303 (2020).

\bibitem{Yildiz2021}
V.~Yildiz,
\newblock \emph{Notes on algebraic structure of truth tables of bracketed formulae connected by implications},
\newblock arXiv:2106.04728 (2021).

\bibitem{FlajoletSedgewick2009}
P.~Flajolet and R.~Sedgewick,
\newblock \emph{Analytic Combinatorics},
\newblock Cambridge University Press, Cambridge, 2009.

\bibitem{Hajek1998}
P.~H{\'a}jek,
\newblock \emph{Metamathematics of Fuzzy Logic},
\newblock Kluwer Academic Publishers, Dordrecht, 1998.



\bibitem{Godel1932}
K.~G{\"o}del,
\newblock \emph{Zum intuitionistischen Aussagenkalk{\"u}l},
\newblock \emph{Anzeiger der Akademie der Wissenschaften in Wien, Mathematisch-naturwissenschaftliche Klasse}
\textbf{69} (1932), 65--66.


\bibitem{CintulaHajekNoguera2011}
P.~Cintula, P.~H{\'a}jek and C.~Noguera (eds.),
\newblock \emph{Handbook of Mathematical Fuzzy Logic, Vol.~1},
\newblock Studies in Logic: Mathematical Logic and Foundations, College Publications, 2011.

\bibitem{OEIS}
N.~J.~A. Sloane et~al.,
\newblock \emph{The On-Line Encyclopedia of Integer Sequences},
\newblock \texttt{https://oeis.org}.

\end{thebibliography}



\vspace{8em}

\begin{verse}
Onlar ki toprakta karınca,\\
suda balık,\\
havada kuş kadar çokturlar;\\
Korkak, cesur, hakim ve çocukturlar.\\
Kahreden ve yaratan ki onlardır,\\
Destanımızda yalnız onların maceraları vardır.
\end{verse}

\begin{flushleft}
Nazım Hikmet.
\end{flushleft}

\end{document}
